% This file is intended to be \input{} from subsubsec__window6_lie_bracket_closure.tex.

\paragraph{根--Cartan 字典与边界塔维数接口的合流：从离散骨架到唯一闭合}
在既定的 window-\(6\) 锚定口径下，“李括号的构造”与“唯一闭合”可被组织为一组可验证的定理化结论，其关键并非引入额外连续结构，而在于把连续李代数所需的输入锚定到离散可审计数据上，并将剩余自由度压缩到有限证书问题。

\begin{remark}[离散 Cartan/根骨架：\texorpdfstring{$18\oplus 3$}{18⊕3} 刚性分裂与根字典输入]\label{rem:window6-rootcartan-18plus3-skeleton}
window-\(6\) 的稳定类型集合满足刚性分解 \(X_6=X_6^{\mathrm{cyc}}\sqcup X_6^{\mathrm{bdry}}\) 且 \(|X_6^{\mathrm{cyc}}|=18\)、\(|X_6^{\mathrm{bdry}}|=3\)。
在根--Cartan 字典口径下，将 boundary 三元组视为 Cartan 锚点、将 cyclic 的 \(18\) 个类型视为根向量标签，则根字典给出一种完全显式的秩 \(3\) 根型输入，并在 \(V=\RR\langle X_6\rangle\) 上诱导紧致 \(B_3/C_3\) 型李括号结构的存在性（定理 \ref{thm:window6-root-dictionary-pullback-bracket}）。
\end{remark}

\begin{remark}[可见规范代数的维数槽位锁定：\texorpdfstring{$(1,3,8)$}{(1,3,8)} \(\Rightarrow\) \texorpdfstring{$\mathfrak u(1)\oplus\mathfrak{su}(2)\oplus\mathfrak{su}(3)$}{u(1)+su(2)+su(3)}]\label{rem:window6-boundary-tower-sm-uniqueness}
边界塔在 \((m=4,6,8)\) 的维数接口给出 \((|X_4^{\mathrm{bdry}}|,|X_6^{\mathrm{bdry}}|,|X_8^{\mathrm{bdry}}|)=(1,3,8)\)。
在 Zeckendorf 原子性与逐项对齐的审计口径下，若可观测连续包络取紧约化，则维数槽位的同构型被唯一锁定为 \(\mathfrak u(1)\oplus\mathfrak{su}(2)\oplus\mathfrak{su}(3)\)（命题 \ref{prop:sm-zeckendorf-lie-algebra-rigidity}）。
\end{remark}

\begin{remark}[高分辨率边界 uplift 的根型分块与对角 \texorpdfstring{$\mathfrak{su}(3)$}{su(3)} 的强迫]\label{rem:window6-m10-a2pm-diagonal-su3}
在 \(m=10\) 的边界 uplift 字典下，根方向的刚性四块分裂把边界层 \(21\) 模态组织为
\[
21\ =\ 6\ (\mathrm{short})\ \oplus\ 6\ (A_2(+))\ \oplus\ 6\ (A_2(-))\ \oplus\ 3\ (\mathrm{Cartan})
\]
且两套 \(A_2(\pm)\) 的判别仅依赖符号结构（定理 \ref{thm:boundary-m10-b3c3-rigid-four-block-split}；表 \ref{tab:boundary_uplift_m10_dictionary}）。
在等权对称被提升为协议不变量并要求其在连续包络上可见时，交换对合的固定点机制排除两份独立的 \(\mathfrak{su}(3)\) 因子在可见层面并存，从而强迫仅保留对角 \(\mathfrak{su}(3)\)（命题 \ref{prop:a2pm-isotropy-diagonal-su3}）。
\end{remark}

上述离散骨架还可进一步凝聚为四条完全结构性的 window-\(6\) 定理：它们分别刻画 cyclic 扇区中 Hamming 权阈值对根长轨道的决定性、\(21=18+3\) 分解对伴随权多重集的精确实现、奇偶荷格中 boundary parity 的零权本质，以及 \(C_6\) 轨道分解与 \(B_3\) 长短根分区之间的唯一交会方式。
\begin{theorem}[window-\(6\) 循环扇区中 Hamming 权阈值对根长 Weyl 轨道的刻画]\label{thm:window6-cyclic-weight-threshold-root-length}
沿用 \S\ref{subsec:bdry-tower-zeck-gut} 的记号，并记
\[
X_6=X_6^{\mathrm{cyc}}\sqcup X_6^{\mathrm{bdry}},
\qquad
|X_6|=21,\quad |X_6^{\mathrm{cyc}}|=18,\quad |X_6^{\mathrm{bdry}}|=3.
\]
cyclic 扇区的 Hamming 权直方图满足
\[
\mathrm{wt}\text{-hist}(X_6^{\mathrm{cyc}})=\{0:1,\ 1:6,\ 2:9,\ 3:2\},
\]
因而恰有六个词满足 \(\mathrm{wt}(w)=1\)，其余十二个词满足 \(\mathrm{wt}(w)\neq 1\)。
则在表 \ref{tab:fold6_b3c3_root_dictionary_b3} 的 \(B_3\) 规范下，\(\mathrm{wt}(w)=1\) 的六个词恰对应短根
\[
\{\pm e_1,\pm e_2,\pm e_3\},
\]
而 \(\mathrm{wt}(w)\neq 1\) 的十二个词恰对应长根
\[
\{\pm e_i\pm e_j:\ 1\le i<j\le 3\}.
\]
在表 \ref{tab:fold6_b3c3_root_dictionary_c3} 的 \(C_3\) 规范下，上述两块仅交换长短根名称：\(\mathrm{wt}(w)=1\) 的六个词恰对应长根
\[
\{\pm 2e_1,\pm 2e_2,\pm 2e_3\},
\]
而 \(\mathrm{wt}(w)\neq 1\) 的十二个词恰对应短根
\[
\{\pm e_i\pm e_j:\ 1\le i<j\le 3\}.
\]
因此 \(X_6^{\mathrm{cyc}}\) 在 \(W(B_3)=W(C_3)\) 的作用下恰分解为两条 Weyl 轨道，轨道大小分别为 \(6\) 与 \(12\)，并由阈值条件 \(\mathrm{wt}=1\) 与 \(\mathrm{wt}\neq 1\) 精确切出。
\end{theorem}
\begin{proof}
表 \ref{tab:fold6_b3c3_root_dictionary_b3} 中 \(\mathrm{wt}=1\) 的六行恰为
\[
\texttt{100000},\ \texttt{010000},\ \texttt{001000},\ \texttt{000100},\ \texttt{000010},\ \texttt{000001},
\]
其像分别为
\[
(1,0,0),\ (0,1,0),\ (0,0,1),\ (0,0,-1),\ (0,-1,0),\ (-1,0,0),
\]
即 \(\{\pm e_1,\pm e_2,\pm e_3\}\)。
另一方面，唯一的 \(\mathrm{wt}=0\) 词 \(\texttt{000000}\) 被送到 \((1,1,0)\)，已经落在
\(\{\pm e_i\pm e_j\}\) 中；其余 \(\mathrm{wt}=2,3\) 的十一行正好给出该集合中的其余十一条向量。
故在 \(B_3\) 规范下，\(\mathrm{wt}=1\) 与 \(\mathrm{wt}\neq 1\) 分别对应短根与长根。

对 \(C_3\) 规范，表 \ref{tab:fold6_b3c3_root_dictionary_c3} 仅将轴向六词统一改写为
\(\{\pm 2e_1,\pm 2e_2,\pm 2e_3\}\)，而其余十二条向量仍为 \(\{\pm e_i\pm e_j\}\)，故结论中的长短根名称恰好交换。

最后，\(W(B_3)=W(C_3)\) 作为 signed permutation 群分别在两类根长集合上传递，因此其拉回到 \(X_6^{\mathrm{cyc}}\) 上正给出大小为 \(6\) 与 \(12\) 的两条 Weyl 轨道。证毕。
\end{proof}

\begin{theorem}[window-\(6\) 的 \texorpdfstring{$21=18+3$}{21=18+3} 分解精确实现 \texorpdfstring{$B_3/C_3$}{B3/C3} 伴随权多重集]\label{thm:window6-21-adjoint-weight-multiset}
在表 \ref{tab:fold6_b3c3_root_dictionary_b3} 与表 \ref{tab:fold6_b3c3_root_dictionary_c3} 的两种规范下，集合 \(X_6\) 都可被规范识别为相应秩 \(3\) 半单 Lie 代数的伴随权多重集：
\[
X_6\ \cong\ \Phi(B_3)\sqcup \{0,0,0\},
\qquad
X_6\ \cong\ \Phi(C_3)\sqcup \{0,0,0\}.
\]
其中 cyclic 扇区 \(X_6^{\mathrm{cyc}}\) 给出全部 \(18\) 条根，而 boundary 三词
\[
X_6^{\mathrm{bdry}}=\{\texttt{100001},\texttt{100101},\texttt{101001}\}
\]
给出三重零权。换言之，window-\(6\) 的 \(21\) 个稳定类型不是“根与外点”的松散并置，而是一个完全显式的秩 \(3\) 伴随权模型。
\end{theorem}
\begin{proof}
两张根字典表分别列出 \(X_6^{\mathrm{cyc}}\) 的 \(18\) 个词及其像，并且像恰为 \(\Phi(B_3)\) 或 \(\Phi(C_3)\) 的全部 \(18\) 条根。
因此 cyclic 扇区在两种规范下都与相应根系完全一致。
另一方面，boundary 扇区恰由三条词
\(\texttt{100001},\texttt{100101},\texttt{101001}\)
组成，且与 cyclic 扇区不交。
于是把这三条 boundary 词规范地视为零权，即得
\[
X_6\cong \Phi(B_3)\sqcup\{0,0,0\}
\qquad\text{或}\qquad
X_6\cong \Phi(C_3)\sqcup\{0,0,0\}.
\]
对秩 \(3\) 的 \(B_3\) 与 \(C_3\)，伴随表示的权多重集正由全部根与三重零权组成，故结论成立。证毕。
\end{proof}

\paragraph{伴随权模型的二阶与四阶不变量}
为把上述 \(18+3\) 权重模型进一步提升为可计算张量不变量，定义两种权映射
\[
\Lambda_B:X_6\longrightarrow \ZZ^3,
\qquad
\Lambda_C:X_6\longrightarrow \ZZ^3
\]
如下：在 cyclic 扇区 \(X_6^{\mathrm{cyc}}\) 上，分别取表 \ref{tab:fold6_b3c3_root_dictionary_b3} 与表 \ref{tab:fold6_b3c3_root_dictionary_c3} 中给出的根向量；在 boundary 三元组 \(X_6^{\mathrm{bdry}}\) 上，统一规定
\[
\Lambda_B(w)=0,
\qquad
\Lambda_C(w)=0.
\]
由定理 \ref{thm:window6-21-adjoint-weight-multiset}，\(\Lambda_B(X_6)\) 与 \(\Lambda_C(X_6)\) 作为多重集分别恰为
\[
\Phi(B_3)\sqcup \{0,0,0\},
\qquad
\Phi(C_3)\sqcup \{0,0,0\}.
\]

\begin{theorem}[window-\(6\) 伴随权模型的二次各向同性与精确 frame 常数]\label{thm:window6-b3c3-adjoint-second-moment-isotropy}
沿用上述 \(\Lambda_B,\Lambda_C\) 的记号。则有精确二阶矩恒等式
\[
\sum_{w\in X_6}\Lambda_B(w)\Lambda_B(w)^{\mathsf T}=10\,I_3,
\qquad
\sum_{w\in X_6}\Lambda_C(w)\Lambda_C(w)^{\mathsf T}=16\,I_3.
\]
因此，\(X_6\) 在两种规范下都给出 \(\RR^3\) 上的各向同性紧框架；其归一化 frame 常数分别为
\[
\frac{10}{3}
\qquad\text{与}\qquad
\frac{16}{3}.
\]
等价地，对任意 \(u\in \RR^3\)，都有
\[
\sum_{w\in X_6}\langle \Lambda_B(w),u\rangle^2=10\,|u|^2,
\qquad
\sum_{w\in X_6}\langle \Lambda_C(w),u\rangle^2=16\,|u|^2.
\]
\end{theorem}
\begin{proof}
boundary 三元组被送到零向量，故不贡献二阶矩；只需对 \(18\) 个根求和。

在 \(B_3\) 规范下，短根部分为
\[
\{\pm e_1,\pm e_2,\pm e_3\},
\]
故
\[
\sum_{\alpha\in \{\pm e_1,\pm e_2,\pm e_3\}}
\alpha\alpha^{\mathsf T}
=2I_3.
\]
长根部分为
\[
\{\pm e_i\pm e_j:\ 1\le i<j\le 3\}.
\]
对固定 \(i<j\)，四个向量
\[
\pm e_i\pm e_j
\]
满足
\[
\sum_{\epsilon_i,\epsilon_j=\pm 1}
(\epsilon_i e_i+\epsilon_j e_j)(\epsilon_i e_i+\epsilon_j e_j)^{\mathsf T}
=4(E_{ii}+E_{jj}),
\]
交叉项由符号对称性抵消。对三对 \((i,j)\) 求和后，每个对角元恰出现两次，故长根总贡献为
\[
8I_3.
\]
于是
\[
\sum_{w\in X_6}\Lambda_B(w)\Lambda_B(w)^{\mathsf T}=2I_3+8I_3=10I_3.
\]

在 \(C_3\) 规范下，长根部分为
\[
\{\pm 2e_1,\pm 2e_2,\pm 2e_3\},
\]
其贡献为
\[
\sum_{i=1}^{3}\bigl((2e_i)(2e_i)^{\mathsf T}+(-2e_i)(-2e_i)^{\mathsf T}\bigr)=8I_3.
\]
其余十二条短根仍为 \(\{\pm e_i\pm e_j\}\)，贡献同上仍为 \(8I_3\)。故
\[
\sum_{w\in X_6}\Lambda_C(w)\Lambda_C(w)^{\mathsf T}=8I_3+8I_3=16I_3.
\]
对任意 \(u\in\RR^3\) 取二次型，即得最后一组等式。证毕。
\end{proof}

\begin{theorem}[window-\(6\) 伴随权模型的指数特征函数与四次 Weyl 不变量]\label{thm:window6-b3c3-adjoint-character-and-quartic-invariant}
沿用 \(\Lambda_B,\Lambda_C\) 的记号。定义指数和
\[
\Theta_B(t):=\sum_{w\in X_6}e^{\langle \Lambda_B(w),t\rangle},
\qquad
\Theta_C(t):=\sum_{w\in X_6}e^{\langle \Lambda_C(w),t\rangle},
\qquad
t=(t_1,t_2,t_3)\in \RR^3.
\]
则有完全显式公式
\[
\Theta_B(t)
=
3+2\sum_{i=1}^{3}\cosh t_i
+4\sum_{1\le i<j\le 3}\cosh t_i\,\cosh t_j,
\]
\[
\Theta_C(t)
=
3+2\sum_{i=1}^{3}\cosh(2t_i)
+4\sum_{1\le i<j\le 3}\cosh t_i\,\cosh t_j.
\]
因此二者分别是 \(B_3\)、\(C_3\) 伴随权多重集在分裂 Cartan 上的精确指数特征函数。

进一步，对任意 \(u=(u_1,u_2,u_3)\in\RR^3\)，令 \(t=su\)。则有四次矩闭式
\[
\sum_{w\in X_6}\langle \Lambda_B(w),u\rangle^4
=
10\sum_{i=1}^{3}u_i^4+24\sum_{1\le i<j\le 3}u_i^2u_j^2,
\]
\[
\sum_{w\in X_6}\langle \Lambda_C(w),u\rangle^4
=
40\sum_{i=1}^{3}u_i^4+24\sum_{1\le i<j\le 3}u_i^2u_j^2.
\]
\end{theorem}
\begin{proof}
由定理 \ref{thm:window6-21-adjoint-weight-multiset}，
\[
\Theta_B(t)=3+\sum_{\alpha\in \Phi(B_3)}e^{\langle \alpha,t\rangle},
\qquad
\Theta_C(t)=3+\sum_{\beta\in \Phi(C_3)}e^{\langle \beta,t\rangle}.
\]
把 \(B_3\) 根系拆成短根 \(\{\pm e_i\}\) 与长根 \(\{\pm e_i\pm e_j\}\)，则
\[
\sum_{\alpha\in \{\pm e_1,\pm e_2,\pm e_3\}}e^{\langle \alpha,t\rangle}
=2\sum_{i=1}^{3}\cosh t_i,
\]
而对每个 \(i<j\)，有
\[
e^{t_i+t_j}+e^{t_i-t_j}+e^{-t_i+t_j}+e^{-t_i-t_j}
=4\cosh t_i\,\cosh t_j.
\]
对三对 \((i,j)\) 求和，即得 \(\Theta_B\) 的闭式。

对 \(C_3\) 规范，仅把轴向六根改为 \(\{\pm 2e_i\}\)，故
\[
\sum_{\beta\in \{\pm 2e_1,\pm 2e_2,\pm 2e_3\}}e^{\langle \beta,t\rangle}
=2\sum_{i=1}^{3}\cosh(2t_i),
\]
而其余十二根仍给出同样的
\(
4\sum_{i<j}\cosh t_i\cosh t_j
\)，遂得 \(\Theta_C\) 的闭式。

现令 \(t=su\)。利用展开式
\[
\cosh(su_i)=1+\frac{s^2u_i^2}{2}+\frac{s^4u_i^4}{24}+O(s^6),
\]
\[
\cosh(2su_i)=1+2s^2u_i^2+\frac{2}{3}s^4u_i^4+O(s^6),
\]
代入上述两条闭式并比较 \(s^4\) 项，可得
\[
\Theta_B(su)
=
21+5s^2|u|^2
+s^4\!\left(
\frac{5}{12}\sum_{i=1}^{3}u_i^4
+\sum_{1\le i<j\le 3}u_i^2u_j^2
\right)+O(s^6),
\]
\[
\Theta_C(su)
=
21+8s^2|u|^2
+s^4\!\left(
\frac{5}{3}\sum_{i=1}^{3}u_i^4
+\sum_{1\le i<j\le 3}u_i^2u_j^2
\right)+O(s^6).
\]
另一方面，由定义
\[
\Theta_B(su)=\sum_{w\in X_6}\exp\!\bigl(s\langle \Lambda_B(w),u\rangle\bigr),
\]
\[
\Theta_C(su)=\sum_{w\in X_6}\exp\!\bigl(s\langle \Lambda_C(w),u\rangle\bigr),
\]
而权集关于原点成对对称，故奇次项消失，且 \(s^4\) 系数分别等于
\[
\frac{1}{24}\sum_{w\in X_6}\langle \Lambda_B(w),u\rangle^4,
\qquad
\frac{1}{24}\sum_{w\in X_6}\langle \Lambda_C(w),u\rangle^4.
\]
将两边比较，即得所述四次矩闭式。证毕。
\end{proof}

\paragraph{三次 Euclidean cubature 与唯一四次调和检测器}
在上述 \(B_3/C_3\) 伴随权模型中，三重 boundary 零权并不只提供静态权多重集分解；它还把 window-\(6\) 的 \(21\) 点配置提升为一对三次精确的 Euclidean cubature，并使四次层面的全部非径向误差收束到唯一一条 \(W(B_3)\)-不变调和方向。

\begin{theorem}[\texorpdfstring{$B_3$}{B3} 归一化下的三次精确 Euclidean cubature]\label{thm:window6-b3-degree3-euclidean-cubature}
\leanverified{paper\_window6\_b3\_degree3\_euclidean\_cubature}
沿用 \(\Lambda_B\) 的记号，并把 \(X_6\) 的像视为带重数的离散配置
\[
\mathcal X_B:=\Lambda_B(X_6)
=
\{0,0,0\}\sqcup \{\pm e_1,\pm e_2,\pm e_3\}\sqcup \{\pm e_i\pm e_j:\ 1\le i<j\le 3\}.
\]
记 \(\sigma_r\) 为半径 \(r\) 的球面概率测度，并定义径向参考测度
\[
\nu_B:=\frac{3}{21}\delta_0+\frac{6}{21}\sigma_1+\frac{12}{21}\sigma_{\sqrt2}.
\]
则对任意总次数不超过 \(3\) 的实多项式 \(P\in\RR[x_1,x_2,x_3]\)，成立严格恒等式
\[
\frac1{21}\sum_{x\in \mathcal X_B}P(x)=\int_{\RR^3}P\,\dd\nu_B.
\]
\end{theorem}
\begin{proof}
只需在总次数不超过 \(3\) 的单项式基上验证。

常数项显然成立。又由于 \(\mathcal X_B\) 关于原点中心对称，任意奇次单项式在离散平均下都为零；而 \(\nu_B\) 是径向测度，因此其奇次矩同样全部为零。故只需检查二次矩。

由定理 \ref{thm:window6-b3c3-adjoint-second-moment-isotropy}，
\[
\sum_{x\in\mathcal X_B}xx^{\mathsf T}=10I_3.
\]
因此
\[
\frac1{21}\sum_{x\in\mathcal X_B}x_ix_j=0
\qquad(i\neq j),
\]
并且
\[
\frac1{21}\sum_{x\in\mathcal X_B}x_i^2=\frac{10}{21}
\qquad(i=1,2,3).
\]
另一方面，由球面对称性，
\[
\int x_ix_j\,\dd\sigma_r=0\qquad(i\neq j),
\qquad
\int x_i^2\,\dd\sigma_r=\frac{r^2}{3}.
\]
于是
\[
\int x_i^2\,\dd\nu_B
=
\frac{6}{21}\cdot\frac13+\frac{12}{21}\cdot\frac23
=
\frac{10}{21},
\]
且交叉矩同样为零。

因此离散平均与 \(\nu_B\) 的全部 \(0,1,2,3\) 阶矩一致，所示 cubature 恒等式成立。证毕。
\end{proof}

\begin{theorem}[\texorpdfstring{$C_3$}{C3} 归一化下的三次精确 Euclidean cubature]\label{thm:window6-c3-degree3-euclidean-cubature}
\leanverified{paper\_window6\_c3\_degree3\_euclidean\_cubature}
沿用 \(\Lambda_C\) 的记号，并把 \(X_6\) 的像视为带重数的离散配置
\[
\mathcal X_C:=\Lambda_C(X_6)
=
\{0,0,0\}\sqcup \{\pm 2e_1,\pm 2e_2,\pm 2e_3\}\sqcup \{\pm e_i\pm e_j:\ 1\le i<j\le 3\}.
\]
定义径向参考测度
\[
\nu_C:=\frac{3}{21}\delta_0+\frac{6}{21}\sigma_{2}+\frac{12}{21}\sigma_{\sqrt2}.
\]
则对任意总次数不超过 \(3\) 的实多项式 \(P\in\RR[x_1,x_2,x_3]\)，成立
\[
\frac1{21}\sum_{x\in \mathcal X_C}P(x)=\int_{\RR^3}P\,\dd\nu_C.
\]
\end{theorem}
\begin{proof}
证明与定理 \ref{thm:window6-b3-degree3-euclidean-cubature} 同样，只需比较二次矩。

由定理 \ref{thm:window6-b3c3-adjoint-second-moment-isotropy}，
\[
\sum_{x\in\mathcal X_C}xx^{\mathsf T}=16I_3.
\]
故
\[
\frac1{21}\sum_{x\in\mathcal X_C}x_ix_j=0
\qquad(i\neq j),
\]
并且
\[
\frac1{21}\sum_{x\in\mathcal X_C}x_i^2=\frac{16}{21}
\qquad(i=1,2,3).
\]
另一方面，
\[
\int x_i^2\,\dd\nu_C
=
\frac{6}{21}\cdot\frac{4}{3}+\frac{12}{21}\cdot\frac{2}{3}
=
\frac{16}{21},
\]
且交叉矩同样因径向对称而为零。

于是离散平均与 \(\nu_C\) 的全部 \(0,1,2,3\) 阶矩一致，结论成立。证毕。
\end{proof}

\begin{theorem}[四次误差的秩一性与唯一的 \texorpdfstring{$W(B_3)$}{W(B3)}-不变调和检测器]\label{thm:window6-b3c3-quartic-rankone-harmonic-detector}
\leanverified{paper\_window6\_b3c3\_quartic\_rankone\_harmonic\_detector}
定义
\[
A:=x_1^4+x_2^4+x_3^4,
\qquad
B:=x_1^2x_2^2+x_1^2x_3^2+x_2^2x_3^2,
\]
\[
H_4:=A-3B,
\qquad
R_4:=(x_1^2+x_2^2+x_3^2)^2=A+2B.
\]
又定义四次误差泛函
\[
\mathcal D_B(P):=\frac1{21}\sum_{x\in\mathcal X_B}P(x)-\int_{\RR^3}P\,\dd\nu_B,
\]
\[
\mathcal D_C(P):=\frac1{21}\sum_{x\in\mathcal X_C}P(x)-\int_{\RR^3}P\,\dd\nu_C.
\]
则：
\begin{enumerate}
  \item \(H_4\) 是 \(W(B_3)\)-不变调和四次多项式；
  \item 在 \(W(B_3)\)-不变四次多项式空间
  \[
  \mathrm{span}\{R_4,H_4\}
  \]
  上，\(\mathcal D_B\) 与 \(\mathcal D_C\) 都是秩一泛函；
  \item 它们满足精确公式
  \[
  \mathcal D_B(R_4)=0,\qquad \mathcal D_B(H_4)=-\frac{2}{7},
  \]
  \[
  \mathcal D_C(R_4)=0,\qquad \mathcal D_C(H_4)=4.
  \]
\end{enumerate}
因此，三次精确性在四次层面只沿唯一一条 \(W(B_3)\)-不变调和方向 \(H_4\) 失效，而纯径向方向 \(R_4\) 仍被完全积分。
\end{theorem}
\begin{proof}
首先，
\[
\Delta A=12(x_1^2+x_2^2+x_3^2),
\qquad
\Delta B=4(x_1^2+x_2^2+x_3^2),
\]
故
\[
\Delta H_4=\Delta(A-3B)=0.
\]
又 \(A,B\) 都是 signed permutation 对称的，故 \(H_4\) 为 \(W(B_3)\)-不变调和四次多项式。

其次，任意 \(W(B_3)\)-不变四次多项式都可唯一写成 \(\alpha A+\beta B\)。等价地，该二维空间亦可由
\[
R_4=A+2B,
\qquad
H_4=A-3B
\]
作为一组基。

现计算 \(\mathcal D_B\)。在 \(\mathcal X_B\) 上，六个短根给出 \(A=1,B=0\)，十二个长根给出 \(A=2,B=1\)，三重原点给出 \(A=B=0\)。因此
\[
\frac1{21}\sum_{x\in\mathcal X_B}A(x)=\frac{6\cdot 1+12\cdot 2}{21}=\frac{10}{7},
\qquad
\frac1{21}\sum_{x\in\mathcal X_B}B(x)=\frac{12}{21}=\frac{4}{7}.
\]
另一方面，对半径为 \(r\) 的球面概率测度有
\[
\int x_i^4\,\dd\sigma_r=\frac{r^4}{5},
\qquad
\int x_i^2x_j^2\,\dd\sigma_r=\frac{r^4}{15}\quad(i\neq j).
\]
故
\[
\int_{\RR^3}A\,\dd\nu_B
=
\frac{6}{21}\cdot\frac{3}{5}
\,+\,
\frac{12}{21}\cdot\frac{12}{5}
=
\frac{54}{35},
\]
\[
\int_{\RR^3}B\,\dd\nu_B
=
\frac{6}{21}\cdot\frac{1}{5}
\,+\,
\frac{12}{21}\cdot\frac{4}{5}
=
\frac{18}{35}.
\]
从而
\[
\mathcal D_B(A)=\frac{10}{7}-\frac{54}{35}=-\frac{4}{35},
\qquad
\mathcal D_B(B)=\frac{4}{7}-\frac{18}{35}=\frac{2}{35}.
\]
于是
\[
\mathcal D_B(R_4)=\mathcal D_B(A+2B)=0,
\qquad
\mathcal D_B(H_4)=\mathcal D_B(A-3B)=-\frac{2}{7}.
\]

再计算 \(\mathcal D_C\)。在 \(\mathcal X_C\) 上，六个轴向长根给出 \(A=16,B=0\)，十二个短根仍给出 \(A=2,B=1\)，故
\[
\frac1{21}\sum_{x\in\mathcal X_C}A(x)=\frac{6\cdot 16+12\cdot 2}{21}=\frac{40}{7},
\qquad
\frac1{21}\sum_{x\in\mathcal X_C}B(x)=\frac{12}{21}=\frac{4}{7}.
\]
而
\[
\int_{\RR^3}A\,\dd\nu_C
=
\frac{6}{21}\cdot\frac{48}{5}
\,+\,
\frac{12}{21}\cdot\frac{12}{5}
=
\frac{144}{35},
\]
\[
\int_{\RR^3}B\,\dd\nu_C
=
\frac{6}{21}\cdot\frac{16}{5}
\,+\,
\frac{12}{21}\cdot\frac{4}{5}
=
\frac{48}{35}.
\]
因此
\[
\mathcal D_C(A)=\frac{40}{7}-\frac{144}{35}=\frac{8}{5},
\qquad
\mathcal D_C(B)=\frac{4}{7}-\frac{48}{35}=-\frac{4}{5}.
\]
遂有
\[
\mathcal D_C(R_4)=\mathcal D_C(A+2B)=0,
\qquad
\mathcal D_C(H_4)=\mathcal D_C(A-3B)=4.
\]

因为 \(\mathcal D_B,\mathcal D_C\) 都在二维空间 \(\mathrm{span}\{R_4,H_4\}\) 上消去 \(R_4\) 而不消去 \(H_4\)，故二者在该空间上的秩都恰为 \(1\)。证毕。
\end{proof}

\begin{corollary}[双归一化四次判别器的唯一性]\label{cor:window6-b3c3-unique-quartic-detector}
\leanverified{paper\_window6\_b3c3\_unique\_quartic\_detector}
在 \(W(B_3)\)-不变多项式代数的次数不超过 \(4\) 的部分中，同时满足
\[
\mathcal D_B(P)=0,
\qquad
\mathcal D_C(P)=0
\]
的多项式空间恰为
\[
\mathrm{span}\bigl\{1,\ x_1^2+x_2^2+x_3^2,\ (x_1^2+x_2^2+x_3^2)^2\bigr\}.
\]
等价地，任一 \(W(B_3)\)-不变四次多项式 \(P\) 都有唯一分解
\[
P=c_4R_4+c_HH_4,
\]
并满足
\[
\mathcal D_B(P)=-\frac{2}{7}c_H,
\qquad
\mathcal D_C(P)=4c_H.
\]
因此，在次数不超过 \(4\) 的 \(W(B_3)\)-不变世界中，唯一能够同时检测两种归一化偏离纯径向积分的方向正是 \(H_4\)。
\end{corollary}
\begin{proof}
由定理 \ref{thm:window6-b3-degree3-euclidean-cubature} 与定理 \ref{thm:window6-c3-degree3-euclidean-cubature}，所有次数不超过 \(3\) 的多项式都已落在 \(\ker \mathcal D_B\cap \ker \mathcal D_C\) 中。

在四次 \(W(B_3)\)-不变空间中，由定理 \ref{thm:window6-b3c3-quartic-rankone-harmonic-detector}，
\[
\mathcal D_B(R_4)=\mathcal D_C(R_4)=0,
\]
而
\[
\mathcal D_B(H_4)\neq 0,
\qquad
\mathcal D_C(H_4)\neq 0.
\]
故
\[
\ker \mathcal D_B\cap \ker \mathcal D_C
\]
在该二维空间上恰等于 \(\mathrm{span}\{R_4\}\)。

再把低次不变量
\[
1,\qquad x_1^2+x_2^2+x_3^2
\]
加入，即得完整交核空间
\[
\mathrm{span}\bigl\{1,\ x_1^2+x_2^2+x_3^2,\ R_4\bigr\}.
\]
最后，将任意四次不变量写成 \(P=c_4R_4+c_HH_4\)，直接代入上一条定理中的数值公式，即得
\[
\mathcal D_B(P)=-\frac{2}{7}c_H,
\qquad
\mathcal D_C(P)=4c_H.
\]
证毕。
\end{proof}

\paragraph{带 boundary 复制方向的五次球面 cubature 与六次首缺陷}
在三次 Euclidean cubature 与四次调和检测器之外，window-\(6\) 的 \(21\) 个标号可见点在单位球面上还携带一条更刚性的五次 cubature 结构。其关键差别在于：上一节处理的是 \(19\) 个几何不同的归一化支撑，而这里保留 boundary 三词作为沿正坐标轴的复制方向，因此必须把 \(21\) 个标号点整体视为一族带重数的球面原子。

为此，沿用表 \ref{tab:fold6_b3c3_root_dictionary_b3} 的 \(B_3\) 字典，并记与 \(\pm e_i\) 对应的 cyclic 短根词为
\[
s_1^+=\texttt{100000},\qquad s_2^+=\texttt{010000},\qquad s_3^+=\texttt{001000},
\]
\[
s_1^-=\texttt{000001},\qquad s_2^-=\texttt{000010},\qquad s_3^-=\texttt{000100}.
\]
又把 boundary 三词按
\[
b_1=\texttt{100001},\qquad b_2=\texttt{100101},\qquad b_3=\texttt{101001}
\]
排序。对每个 \(w\in X_6\) 定义单位可见向量
\[
u_w:=\frac{v_w}{|v_w|}\in S^2,
\]
其中 \(v_w=\Lambda_B(w)\) 若 \(w\in X_6^{\mathrm{cyc}}\)，而 \(v_{b_i}=e_i\)。
于是 \(\widetilde U_6:=\{u_w:\ w\in X_6\}\) 是一个带三重重复方向
\[
\{e_1,e_2,e_3\}
\]
的 \(21\) 点标号球面支撑。以下记 \(\sigma\) 为 \(S^2\) 上的球面概率测度，并称有限加权原子测度
\[
\mu=\sum_{w\in X_6}c_w\delta_{u_w}
\]
在 \(\widetilde U_6\) 上对次数不超过 \(t\) 的球面多项式给出精确 cubature，若对任意次数 \(\le t\) 的多项式 \(P\) 都有
\[
\sum_{w\in X_6}c_w P(u_w)=\mu(S^2)\int_{S^2}P(u)\,\dd\sigma(u).
\]

\begin{theorem}[window-\(6\) 标号球面支撑上的五次 cubature 族]\label{thm:window6-labeled-spherical-degree5-cubature-family}
设
\[
\mu=\sum_{w\in X_6}c_w\delta_{u_w}
\]
是支撑于 \(\widetilde U_6\) 的有限加权原子测度。则 \(\mu\) 对 \(S^2\) 的全部次数 \(\le 5\) 多项式给出精确球面 cubature，当且仅当存在参数
\[
\lambda,t_1,t_2,t_3\in \CC
\]
使得
\[
c_w=\lambda
\qquad
\text{对全部 \(12\) 个长根方向},
\]
\[
c_{s_i^-}=\frac{\lambda}{2},
\qquad
c_{s_i^+}=\frac{\lambda}{2}-t_i,
\qquad
c_{b_i}=t_i
\qquad (i=1,2,3).
\]
特别地，总质量恒为
\[
\mu(S^2)=15\lambda.
\]
因此，window-\(6\) 的 \(21\) 点标号球面支撑上的五次 cubature 并非离散孤点，而是一个四维仿射族；其真正模只来自三条正轴复制方向上的质量转移。
\end{theorem}
\begin{proof}
对每个 \(i\) 记
\[
m_i^+:=c_{s_i^+}+c_{b_i},
\qquad
m_i^-:=c_{s_i^-}.
\]
由于 \(u_{s_i^+}=u_{b_i}=e_i\)，球面矩条件只依赖于 \(m_i^+\)，而不依赖于它在 \((s_i^+,b_i)\) 之间如何分拆。
因此可先在 \(18\) 个几何不同的方向上分类所有五次 cubature，再把 \(m_i^+\) 在 \((s_i^+,b_i)\) 间自由分配。

对每个坐标平面 \(ij\in\{12,13,23\}\)，记四个长根方向
\[
\frac{\pm e_i\pm e_j}{\sqrt2}
\]
上的权为
\[
\lambda_{ij}^{++},\qquad
\lambda_{ij}^{+-},\qquad
\lambda_{ij}^{-+},\qquad
\lambda_{ij}^{--}.
\]
由于球面概率测度的一切奇次矩都为零，取混合矩
\[
x_i x_j,\qquad x_i x_j^2,\qquad x_i^2 x_j
\]
便得到
\[
\lambda_{ij}^{++}-\lambda_{ij}^{+-}-\lambda_{ij}^{-+}+\lambda_{ij}^{--}=0,
\]
\[
\lambda_{ij}^{++}+\lambda_{ij}^{+-}-\lambda_{ij}^{-+}-\lambda_{ij}^{--}=0,
\]
\[
\lambda_{ij}^{++}-\lambda_{ij}^{+-}+\lambda_{ij}^{-+}-\lambda_{ij}^{--}=0.
\]
解此线性系统得
\[
\lambda_{ij}^{++}=\lambda_{ij}^{+-}=\lambda_{ij}^{-+}=\lambda_{ij}^{--}=:\lambda_{ij}.
\]
故每个坐标平面中的四个长根必须等权。

再看一阶矩。由于三组长根在每个坐标方向上的总贡献已成对相消，取线性函数 \(x_i\) 即得
\[
m_i^+=m_i^-
\qquad (i=1,2,3).
\]

最后使用四阶 isotropy。球面概率测度满足
\[
\int_{S^2}x_i^4\,\dd\sigma=3\int_{S^2}x_i^2x_j^2\,\dd\sigma
\qquad(i\ne j).
\]
而在当前离散测度下，
\[
M_{iijj}:=\sum_{w\in X_6}c_w u_{w,i}^2u_{w,j}^2=\lambda_{ij},
\]
\[
M_{iiii}:=\sum_{w\in X_6}c_w u_{w,i}^4
=m_i^++m_i^-+\lambda_{ij}+\lambda_{ik}
=2m_i^++\lambda_{ij}+\lambda_{ik}.
\]
因此
\[
2m_i^+ + \lambda_{ij}+\lambda_{ik}=3\lambda_{ij}=3\lambda_{ik}.
\]
由此推出
\[
\lambda_{12}=\lambda_{13}=\lambda_{23}=:\lambda,
\qquad
m_i^+=m_i^-=\frac{\lambda}{2}.
\]
于是全部长根权都等于 \(\lambda\)，全部负短根权都等于 \(\lambda/2\)，而每条正轴上的总质量都等于 \(\lambda/2\)。
把该总质量在 \((s_i^+,b_i)\) 间分拆为
\[
c_{s_i^+}=\frac{\lambda}{2}-t_i,
\qquad
c_{b_i}=t_i,
\]
便得到所述参数化。总质量随即为
\[
12\lambda+3\cdot\frac{\lambda}{2}+3\cdot\frac{\lambda}{2}=15\lambda.
\]

反过来，若权满足题设公式，则一切奇次矩都因长根的符号对称与轴向的平衡条件
\[
m_i^+=m_i^-=\frac{\lambda}{2}
\]
而消失。对二次矩，三个坐标轴方向总贡献为 \(\lambda I_3\)，三组长根方向总贡献为 \(4\lambda I_3\)，故
\[
\sum_{w\in X_6}c_w\langle x,u_w\rangle^2=5\lambda |x|^2.
\]
对四次矩，三个坐标轴方向总贡献为
\[
\lambda(x_1^4+x_2^4+x_3^4),
\]
而每个坐标平面的四个长根方向给出
\[
\lambda\bigl(x_i^4+6x_i^2x_j^2+x_j^4\bigr).
\]
对三对 \((i,j)\) 求和后得到
\[
\sum_{w\in X_6}c_w\langle x,u_w\rangle^4=3\lambda |x|^4.
\]
因此
\[
\sum_{w\in X_6}c_w\langle x,u_w\rangle^k
=
15\lambda\int_{S^2}\langle x,u\rangle^k\,\dd\sigma(u)
\qquad (k=0,1,2,3,4,5),
\]
其中用到了
\[
\int_{S^2}\langle x,u\rangle^2\,\dd\sigma(u)=\frac{|x|^2}{3},
\qquad
\int_{S^2}\langle x,u\rangle^4\,\dd\sigma(u)=\frac{|x|^4}{5},
\]
以及奇次矩为零。由极化恒等式，这等价于对全部次数 \(\le 5\) 多项式的精确 cubature。证毕。
\end{proof}
\leanverified{paper\_window6\_labeled\_spherical\_degree5\_cubature\_family}

\begin{proposition}[正权归一化五次 cubature 的立方体模空间]\label{prop:window6-labeled-spherical-degree5-cubature-cube}
在定理 \ref{thm:window6-labeled-spherical-degree5-cubature-family} 的四维仿射族中，正权且归一化
\[
\sum_{w\in X_6}c_w=1
\]
的部分恰为闭立方体
\[
\lambda=\frac{1}{15},
\qquad
0\le t_i\le \frac{1}{30}
\qquad (i=1,2,3).
\]
故 window-\(6\) 的正权五次球面 cubature 模空间同构于
\[
{\Bigl[0,\frac{1}{30}\Bigr]}^{3}.
\]
其 \(8\) 个顶点给出全部支撑极小的正权解；在这些顶点上，支撑恰有 \(18\) 个点，并且顶点集天然构成一个 \(\ZZ_2^3\)-torsor。

此外，在整个 \(21\) 点标号支撑上不存在非平凡等权五次球面设计。
\end{proposition}
\begin{proof}
由定理 \ref{thm:window6-labeled-spherical-degree5-cubature-family}，总质量恒为 \(15\lambda\)，故归一化立即给出
\[
\lambda=\frac{1}{15}.
\]
正权条件等价于
\[
\lambda>0,
\qquad
\frac{\lambda}{2}-t_i\ge 0,
\qquad
t_i\ge 0.
\]
代入 \(\lambda=1/15\) 即得
\[
0\le t_i\le \frac{1}{30}.
\]
故正权归一化模空间正是三维闭立方体。

在该立方体中，\(12\) 个长根权恒为 \(1/15\)，三条负短根权恒为 \(1/30\)，而每条正轴上的总质量 \(1/30\) 在 \((s_i^+,b_i)\) 间分配。若 \(t_i\) 处于开区间，则 \(s_i^+\) 与 \(b_i\) 两个标号点都带正权；若 \(t_i\in\{0,1/30\}\)，则恰有一个点的权消失。故支撑最小要求三个 \(t_i\) 同时落在端点，从而一共保留
\[
12+3+3=18
\]
个非零权点。三条正轴各有两种选择，因而顶点集与三维超立方体的顶点集自然同构，即构成一个 \(\ZZ_2^3\)-torsor。

最后，若全部 \(21\) 个标号点都带同一非零权 \(c\)，则定理 \ref{thm:window6-labeled-spherical-degree5-cubature-family} 强迫
\[
c=\lambda
\qquad\text{在全部长根方向上},
\]
同时又有
\[
c=\frac{\lambda}{2}
\qquad\text{在三条负短根上},
\]
矛盾。故非平凡等权五次球面设计不存在。证毕。
\end{proof}

\begin{theorem}[整个五次 cubature 立方体共享同一六次谐和缺陷]\label{thm:window6-labeled-spherical-degree6-harmonic-defect}
设 \(\mu\) 是定理 \ref{thm:window6-labeled-spherical-degree5-cubature-family} 的任意五次精确 cubature。定义
\[
A_6(x):=x_1^6+x_2^6+x_3^6,
\qquad
B_6(x):=\sum_{i\ne j}x_i^4x_j^2,
\qquad
C_6(x):=x_1^2x_2^2x_3^2,
\]
以及六次谐和多项式
\[
\mathscr H_6(x):=
2A_6(x)-15B_6(x)+180C_6(x).
\]
则对任意 \(x\in \RR^3\)，有精确恒等式
\[
\sum_{w\in X_6}c_w\langle x,u_w\rangle^6
=
\frac{15\lambda}{7}|x|^6-\frac{\lambda}{14}\mathscr H_6(x),
\]
其中 \(\lambda\) 为定理 \ref{thm:window6-labeled-spherical-degree5-cubature-family} 中的长根公共权。
在归一化情形 \(\lambda=1/15\) 下，公式化为
\[
\sum_{w\in X_6}c_w\langle x,u_w\rangle^6
=
\frac{1}{7}|x|^6-\frac{1}{210}\mathscr H_6(x).
\]
特别地，六次矩缺陷与参数 \((t_1,t_2,t_3)\) 无关；整个五次 cubature 立方体在六次上共享同一个首个不可消去 defect。
\end{theorem}
\begin{proof}
由定理 \ref{thm:window6-labeled-spherical-degree5-cubature-family}，三条正轴上的总质量始终满足
\[
c_{s_i^+}+c_{b_i}=\frac{\lambda}{2},
\]
故六次矩只依赖于 \(\lambda\)，与 \((t_1,t_2,t_3)\) 如何分配无关。

沿三个坐标轴的总贡献为
\[
\lambda A_6(x).
\]
另一方面，对固定的 \(i<j\)，四个长根方向给出
\[
\lambda\sum_{\epsilon_i,\epsilon_j=\pm 1}
\left\langle x,\frac{\epsilon_i e_i+\epsilon_j e_j}{\sqrt2}\right\rangle^{6}
=
\frac{\lambda}{8}
\sum_{\epsilon_i,\epsilon_j=\pm 1}
{\bigl(\epsilon_i x_i+\epsilon_j x_j\bigr)}^{6}
=
\frac{\lambda}{2}\bigl(x_i^6+15x_i^4x_j^2+15x_i^2x_j^4+x_j^6\bigr).
\]
对三对 \((i,j)\) 求和后得到
\[
\sum_{w\in X_6}c_w\langle x,u_w\rangle^6
=
2\lambda A_6(x)+\frac{15\lambda}{2}B_6(x).
\]

又由
\[
|x|^6=A_6(x)+3B_6(x)+6C_6(x)
\]
可知
\[
\frac{15\lambda}{7}|x|^6-\frac{\lambda}{14}\mathscr H_6(x)
=
\lambda\left(
\frac{15}{7}(A_6+3B_6+6C_6)-\frac{1}{14}(2A_6-15B_6+180C_6)
\right),
\]
整理即得
\[
\frac{15\lambda}{7}|x|^6-\frac{\lambda}{14}\mathscr H_6(x)
=
2\lambda A_6(x)+\frac{15\lambda}{2}B_6(x).
\]
这就证明了题设恒等式。

最后，直接计算拉普拉斯算子可得
\[
\Delta \mathscr H_6=0,
\]
故 \(\mathscr H_6\) 的确是六次谐和多项式。归一化情形只需代入 \(\lambda=1/15\)。证毕。
\end{proof}

\begin{corollary}[window-\(6\) 标号球面支撑的 cubature 强度恰为 \(5\)]\label{cor:window6-labeled-spherical-cubature-strength-five}
window-\(6\) 的 \(21\) 点标号球面支撑 \(\widetilde U_6\) 上存在非平凡的五次精确球面 cubature；但不存在任何总质量非零、支撑于同一 \(21\) 点标号支撑的六次精确球面 cubature。
因此其最大精确阶恰为
\[
t_{\max}=5.
\]
\end{corollary}
\begin{proof}
存在性已由定理 \ref{thm:window6-labeled-spherical-degree5-cubature-family} 给出。
为证六次不可能，只需取球面六次单项式
\[
P(x):=x_1^2x_2^2x_3^2.
\]
对每个 \(w\in X_6\)，向量 \(u_w\) 至少有一个坐标为零：长根方向只占据两个坐标，而短根与 boundary 方向只占据一个坐标。故
\[
P(u_w)=0
\qquad(\forall\,w\in X_6),
\]
从而任意支撑于 \(\widetilde U_6\) 的原子测度 \(\mu\) 都满足
\[
\sum_{w\in X_6}c_w P(u_w)=0.
\]
另一方面，
\[
\int_{S^2}x_1^2x_2^2x_3^2\,\dd\sigma=\frac{1}{105}>0.
\]
因此若 \(\mu(S^2)\ne 0\)，便不可能有
\[
\sum_{w\in X_6}c_w P(u_w)=\mu(S^2)\int_{S^2}P(u)\,\dd\sigma(u).
\]
故六次精确 cubature 不存在。证毕。
\end{proof}

\begin{conjecture}[boundary 复制方向的低阶不可见性与首个统一谐和缺陷]\label{conj:window-even-boundary-duplicate-spherical-cubature}
设 \(m\) 为偶数，并假定该窗口存在与 window-\(6\) 同型的可见分裂
\[
X_m=X_m^{\mathrm{cyc}}\sqcup X_m^{\mathrm{bdry}},
\]
其中 \(X_m^{\mathrm{cyc}}\) 经某个审计字典嵌入约化根系
\[
\Phi_m\subset \mathfrak h_m^\ast,
\]
而 \(X_m^{\mathrm{bdry}}\) 恰给出若干与正短根方向重复的 boundary copies。
则应当存在以下普适现象：
\begin{enumerate}
  \item 低阶 \(SO(r_m)\)-moment exactness 只依赖于每条重复方向上的总质量，而不依赖于该总质量在 cyclic copy 与 boundary copy 之间的分配；
  \item 正权精确 cubature 的归一化模空间是一个 \(r_m\)-维立方体；
  \item 第一个超出精确阶的缺陷由一个与立方体参数无关的 \(W(\Phi_m)\)-谐和多项式统一控制。
\end{enumerate}
window-\(6\) 正是该原则的首个经审计样本：此时 \(r_6=3\)，正权模空间为
\[
{\Bigl[0,\frac{1}{30}\Bigr]}^{3},
\]
而首个缺陷由定理 \ref{thm:window6-labeled-spherical-degree6-harmonic-defect} 中的 \(\mathscr H_6\) 唯一控制。
\end{conjecture}

\paragraph{几何解释}
定理 \ref{thm:window6-labeled-spherical-degree5-cubature-family}--\ref{thm:window6-labeled-spherical-degree6-harmonic-defect} 表明：在一切不超过五阶的 \(SO(3)\) 角分辨可观测量中，boundary 正轴质量向 cyclic 正轴质量的转移完全不可见；真正开始分辨两者差异的，不是一阶偶极，也不是四阶各向异性，而是同一个六次谐和 defect \(\mathscr H_6\)。因此，window-\(6\) 的 boundary 扇区并不是低阶运动学中的显性扰动，而是高阶角动量通道中的首个可测破缺源。

\paragraph{根系支撑的精确消失理想、Hilbert twin 现象与四次插值饱和}
在 \(B_3/C_3\) 根字典、三次 Euclidean cubature 与唯一四次调和检测器都已固定之后，还可把 window-\(6\) 的 \(19\) 点归一化支撑进一步提升为一个完全显式的零维代数几何对象：其消失理想、标准单项式基、Hilbert 系列以及四次以内的全部插值关系都可被闭式写出。

\begin{theorem}[\texorpdfstring{$B_3$}{B3} 归一化支撑的精确消失理想与 Hilbert 系列]\label{thm:window6-b3-support-vanishing-ideal-hilbert}
沿用定理 \ref{thm:window6-21-adjoint-weight-multiset} 中的 \(B_3\) 规范，并记
\[
\mathscr X_B:=\{0\}\cup \Phi(B_3)\subset \mathbb A^3.
\]
则其消失理想恰为
\[
I_B=
\bigl\langle
x_1^3-x_1,\ x_2^3-x_2,\ x_3^3-x_3,\ x_1x_2x_3
\bigr\rangle
\subset \mathbb Q[x_1,x_2,x_3].
\]
并且在词典序 \(x_1>x_2>x_3\) 下，标准单项式恰为
\[
\mathcal B_B=
\bigl\{
x_1^a x_2^b x_3^c:\ 0\le a,b,c\le 2,\ \min(a,b,c)=0
\bigr\}.
\]
从而
\[
\dim_{\mathbb Q}\mathbb Q[x_1,x_2,x_3]/I_B=19,
\]
且其 Hilbert 系列为
\[
\boxed{\;
H_B(t)=1+3t+6t^2+6t^3+3t^4.\;}
\]
\leanpartial{paper\_window6\_b3\_support\_count}{仅形式化 19-点计数代数核 $|\{x \in \{-1,0,1\}^3 : x_1 x_2 x_3 = 0\}| = 27 - 8 = 19$；消失理想 $I_B$、Hilbert 系列 $H_B(t)$ 与 $B_3$ 根系结构留后轮}
\leanverified{cube27\_card}
\leanverified{supportNonzero\_card}
\leanverified{supportZero\_disjoint\_nonzero}
\leanverified{supportZero\_union\_nonzero}
\leanverified{supportZero\_card}
\end{theorem}
\begin{proof}
由方程
\[
x_i^3-x_i=0
\qquad (i=1,2,3)
\]
可知任一零点都满足
\[
x_i\in\{-1,0,1\}.
\]
再加上约束
\[
x_1x_2x_3=0,
\]
便得到
\[
V(I_B)
=
\bigl\{
(x_1,x_2,x_3)\in\{-1,0,1\}^3:\ x_1x_2x_3=0
\bigr\}.
\]
其基数为
\[
3^3-2^3=19.
\]
另一方面，\(\Phi(B_3)\) 在标准模型下恰为
\[
\{\pm e_1,\pm e_2,\pm e_3\}
\sqcup
\{\pm e_i\pm e_j:\ 1\le i<j\le 3\},
\]
亦即所有坐标属于 \(\{-1,0,1\}\)、且至少一个坐标为 \(0\) 的非零点。因此
\[
\mathscr X_B=\{0\}\cup \Phi(B_3)=V(I_B).
\]

现取词典序 \(x_1>x_2>x_3\)。四个生成元的首项分别为
\[
x_1^3,\qquad x_2^3,\qquad x_3^3,\qquad x_1x_2x_3.
\]
由于前三个首项两两互素，而其余与 \(x_1x_2x_3\) 的最小公倍式对应的 \(S\)-多项式都可直接由同一生成系约化到 \(0\)，故该生成系本身即为一个 Gröbner 基。于是
\[
\mathrm{in}(I_B)=\langle x_1^3,x_2^3,x_3^3,x_1x_2x_3\rangle.
\]
因此标准单项式正是所有指数都不超过 \(2\) 且不同时含三变量的单项式，即
\[
\mathcal B_B=
\bigl\{
x_1^a x_2^b x_3^c:\ 0\le a,b,c\le 2,\ \min(a,b,c)=0
\bigr\}.
\]
其个数为
\[
3^3-2^3=19.
\]
故
\[
\dim_{\mathbb Q}\mathbb Q[x_1,x_2,x_3]/I_B=19.
\]

又因 \(I_B\subseteq I(\mathscr X_B)\)，故存在满射
\[
\mathbb Q[x_1,x_2,x_3]/I_B\twoheadrightarrow \mathbb Q[x_1,x_2,x_3]/I(\mathscr X_B),
\]
右端维数等于 \(|\mathscr X_B|=19\)。左端维数亦为 \(19\)，故该满射只能是同构，从而
\[
I_B=I(\mathscr X_B).
\]

最后按总次数计数标准单项式即可得到
\[
H_B(t)=1+3t+6t^2+6t^3+3t^4.
\]
证毕。
\end{proof}
\leanverified{paper\_window6\_b3\_support\_vanishing\_ideal\_hilbert}

\begin{theorem}[\texorpdfstring{$C_3$}{C3} 归一化支撑的精确消失理想与 Hilbert 系列]\label{thm:window6-c3-support-vanishing-ideal-hilbert}
沿用定理 \ref{thm:window6-21-adjoint-weight-multiset} 中的 \(C_3\) 规范，并记
\[
\mathscr X_C:=\{0\}\cup \Phi(C_3)\subset \mathbb A^3,
\qquad
R:=x_1^2+x_2^2+x_3^2.
\]
则 \(\mathscr X_C\) 的消失理想可写为
\[
I_C=
\Bigl\langle
x_1(x_1^2-1)(x_1^2-4),\,
x_2(x_2^2-1)(x_2^2-4),\,
x_3(x_3^2-1)(x_3^2-4),\,
x_1x_2x_3,\,
R(R-2)(R-4)
\Bigr\rangle.
\]
并且在词典序 \(x_1>x_2>x_3\) 下，\(\mathbb Q[x_1,x_2,x_3]/I_C\) 的标准单项式恰为
\[
\mathcal B_C=
\{
1,\ x_1,\ x_2,\ x_3,\ x_1^2,\ x_1x_2,\ x_1x_3,\ x_2^2,\ x_2x_3,\ x_3^2,
\]
\[
x_1x_2^2,\ x_1x_3^2,\ x_2^3,\ x_2^2x_3,\ x_2x_3^2,\ x_3^3,\ x_2^4,\ x_2^2x_3^2,\ x_3^4
\}.
\]
从而
\[
\dim_{\mathbb Q}\mathbb Q[x_1,x_2,x_3]/I_C=19,
\]
且其 Hilbert 系列同样为
\[
\boxed{\;
H_C(t)=1+3t+6t^2+6t^3+3t^4.\;}
\]
\leanverified{paper\_window6\_c3\_support\_count}
\end{theorem}
\begin{proof}
前三类生成元逐坐标强制
\[
x_i\in\{0,\pm 1,\pm 2\},
\qquad i=1,2,3.
\]
再由
\[
x_1x_2x_3=0
\]
知至多有两个坐标非零；最后
\[
R(R-2)(R-4)=0
\]
强制平方范数满足
\[
R=x_1^2+x_2^2+x_3^2\in\{0,2,4\}.
\]
因此零点只能分成三类：
\[
(0,0,0),
\]
\[
(\pm 2,0,0),\ (0,\pm 2,0),\ (0,0,\pm 2),
\]
以及
\[
(\pm 1,\pm 1,0),\ (\pm 1,0,\pm 1),\ (0,\pm 1,\pm 1).
\]
它们的总数为
\[
1+6+12=19.
\]
而 \(\Phi(C_3)\) 的标准模型恰为
\[
\{\pm 2e_1,\pm 2e_2,\pm 2e_3\}
\sqcup
\{\pm e_i\pm e_j:\ 1\le i<j\le 3\},
\]
故
\[
V(I_C)=\mathscr X_C.
\]

对该生成系作 Buchberger 消元，在词典序 \(x_1>x_2>x_3\) 下得到的初始单项式理想由
\[
x_1^3,\ x_1^2x_2,\ x_1^2x_3,\ x_1x_2^3,\ x_1x_2x_3,\ x_1x_3^3,\ x_2^5,\ x_2^3x_3,\ x_2x_3^3,\ x_3^5
\]
生成。于是标准单项式恰为未被上述单项式整除的单项式，也就是题述的 \(19\) 个单项式 \(\mathcal B_C\)。因此
\[
\dim_{\mathbb Q}\mathbb Q[x_1,x_2,x_3]/I_C=19.
\]

同样地，由包含 \(I_C\subseteq I(\mathscr X_C)\) 诱导的满射
\[
\mathbb Q[x_1,x_2,x_3]/I_C\twoheadrightarrow \mathbb Q[x_1,x_2,x_3]/I(\mathscr X_C)
\]
以及两端同为 \(19\) 维，可知该满射为同构，从而
\[
I_C=I(\mathscr X_C).
\]
最后按总次数统计 \(\mathcal B_C\) 中单项式，即得
\[
H_C(t)=1+3t+6t^2+6t^3+3t^4.
\]
证毕。
\end{proof}
\leanverified{paper\_window6\_c3\_support\_vanishing\_ideal\_hilbert}

\begin{corollary}[\texorpdfstring{$B_3/C_3$}{B3/C3} 两支撑的 Hilbert twin 现象与四次唯一插值]\label{cor:window6-b3c3-hilbert-twins-degree4-interpolation}
沿用定理 \ref{thm:window6-b3-support-vanishing-ideal-hilbert} 与定理 \ref{thm:window6-c3-support-vanishing-ideal-hilbert} 的记号。则尽管 \(\mathscr X_B\) 与 \(\mathscr X_C\) 在几何上并不等距、其半径分层也不同，二者却满足完全相同的 Hilbert 系列
\[
H_B(t)=H_C(t)=1+3t+6t^2+6t^3+3t^4.
\]
因此它们在标准分次意义下构成一对 Hilbert twins。进一步：
\begin{enumerate}
  \item 对任意函数 \(f:\mathscr X_B\to \mathbb Q\)，存在且仅存在一个多项式
  \[
  P_B\in \mathrm{span}(\mathcal B_B)
  \]
  使得
  \[
  P_B|_{\mathscr X_B}=f;
  \]
  \item 对任意函数 \(g:\mathscr X_C\to \mathbb Q\)，存在且仅存在一个多项式
  \[
  P_C\in \mathrm{span}(\mathcal B_C)
  \]
  使得
  \[
  P_C|_{\mathscr X_C}=g.
  \]
\end{enumerate}
换言之，二者的插值复杂度严格相同，并且都在总次数 \(4\) 处闭合。
\leanverified{paper\_window6\_b3c3\_hilbert\_twins\_degree4\_interpolation}
\end{corollary}
\begin{proof}
第一条等式只是前两条定理的直接并列。

又由定理 \ref{thm:window6-b3-support-vanishing-ideal-hilbert}，
\(\mathcal B_B\) 是坐标环
\[
\mathbb Q[x_1,x_2,x_3]/I_B
\]
的一组标准单项式基，且其中每个单项式的总次数都不超过 \(4\)。因此
\[
\dim \mathrm{span}(\mathcal B_B)=19=|\mathscr X_B|.
\]
由于 \(I_B=I(\mathscr X_B)\)，评价映射
\[
\mathrm{ev}_B:\mathrm{span}(\mathcal B_B)\longrightarrow \mathrm{Fun}(\mathscr X_B,\mathbb Q)
\]
正是坐标环到函数代数的同构在这组基上的实现，故为同构，从而唯一插值成立。

对 \(\mathscr X_C\) 与 \(\mathcal B_C\) 完全同理。证毕。
\end{proof}

\begin{theorem}[\texorpdfstring{$B_3/C_3$}{B3/C3} 支撑的四次正交补空间维数为 \(16\)，且由消失理想完全控制]\label{thm:window6-b3c3-degree4-relation-space-saturation}
记
\[
\mathcal P_{\le 4}:=
\{P\in \mathbb Q[x_1,x_2,x_3]:\deg P\le 4\},
\qquad
\dim \mathcal P_{\le 4}=\binom{4+3}{3}=35.
\]
对 \(S\in\{B,C\}\)，定义评价映射
\[
\mathrm{ev}_S:\mathcal P_{\le 4}\to \mathrm{Fun}(\mathscr X_S,\mathbb Q).
\]
则有
\[
\mathrm{rank}(\mathrm{ev}_S)=19,
\qquad
\dim\ker(\mathrm{ev}_S)=16.
\]
并且
\[
\ker(\mathrm{ev}_B)=I_B\cap \mathcal P_{\le 4},
\qquad
\ker(\mathrm{ev}_C)=I_C\cap \mathcal P_{\le 4}.
\]
特别地，\(\mathscr X_B\) 与 \(\mathscr X_C\) 在四次以内的全部代数关系都已饱和：不存在任何额外的低次数关系超出各自消失理想所给出的那些。
\leanverified{paper\_window6\_b3c3\_degree4\_relation\_space\_saturation}
\end{theorem}
\begin{proof}
由推论 \ref{cor:window6-b3c3-hilbert-twins-degree4-interpolation}，对 \(S=B,C\) 都可在 \(\mathcal P_{\le 4}\) 中找到一个 \(19\) 维子空间，其评价映射到 \(\mathrm{Fun}(\mathscr X_S,\mathbb Q)\) 为同构。因此
\[
\mathrm{rank}(\mathrm{ev}_S)=19.
\]
再由秩--零度定理，
\[
\dim\ker(\mathrm{ev}_S)=35-19=16.
\]

另一方面，\(\ker(\mathrm{ev}_S)\) 按定义就是所有在 \(\mathscr X_S\) 上消失的次数不超过 \(4\) 的多项式，故恰为
\[
I(\mathscr X_S)\cap \mathcal P_{\le 4}.
\]
由定理 \ref{thm:window6-b3-support-vanishing-ideal-hilbert} 与定理 \ref{thm:window6-c3-support-vanishing-ideal-hilbert}，
\[
I(\mathscr X_B)=I_B,
\qquad
I(\mathscr X_C)=I_C,
\]
于是
\[
\ker(\mathrm{ev}_B)=I_B\cap \mathcal P_{\le 4},
\qquad
\ker(\mathrm{ev}_C)=I_C\cap \mathcal P_{\le 4}.
\]
这正说明四次以内的一切关系都已经由各自的完整消失理想给出，而不存在额外独立的低次数约束。证毕。
\end{proof}

\noindent
由此，window-\(6\) 的 \(B_3/C_3\) 根系字典不再只是一个离散标号装置，而是同时诱导出两套 \(19\) 点零维支撑的完整代数几何模型：一方面，二者分别具有显式消失理想与标准单项式基；另一方面，尽管其几何半径结构不同，却共享同一 Hilbert 系列、同一四次插值阈值以及同一 \(16\) 维低次关系空间规模。于是，\(B_3\) 与 \(C_3\) 两种归一化不只在根长交换意义下对偶，而且在有限支撑的分次代数复杂度上也形成一对严格的 Hilbert twins。

\endinput

\paragraph{三张 Levi 平面、三轴选择律与唯一 quartic 残差}
在 \(B_3/C_3\) 可见权字典、三次 cubature、四次调和检测器与支撑理想都已固定之后，window-\(6\) 的可见部门还可进一步压缩为一条更接近有效场论的局域结构链：全部非零支撑并不充满三维 Cartan 空间，而是精确贴在三张 rank-\(2\) Levi 平面上；任何真正的三轴单点耦合都逐点湮灭；二阶之后的首个非标量局域残差只沿唯一一条八面体型四次调和方向存活；而 \(C_3\) 规范相对于 \(B_3\) 规范的全部增厚则可在全偶矩层级上写成闭式。

\begin{theorem}[window-\(6\) 可见权的三 Levi 平面分解]\label{thm:window6-b3c3-visible-support-three-levi-planes}
沿用定理 \ref{thm:window6-21-adjoint-weight-multiset} 中的权映射
\[
\Lambda_B,\Lambda_C:X_6\to \ZZ^3
\]
及 boundary 三词映到零权的约定。
对 \(1\le i<j\le 3\)，记
\[
H_{ij}:=\operatorname{span}\{e_i,e_j\}\subset \RR^3.
\]
则有精确并分解
\[
\Lambda_B(X_6)\setminus\{0\}
=
\Phi^{B_2}_{12}\cup \Phi^{B_2}_{13}\cup \Phi^{B_2}_{23},
\qquad
\Phi^{B_2}_{ij}:=\{\pm e_i,\pm e_j,\pm e_i\pm e_j\}\subset H_{ij},
\]
以及
\[
\Lambda_C(X_6)\setminus\{0\}
=
\Phi^{C_2}_{12}\cup \Phi^{C_2}_{13}\cup \Phi^{C_2}_{23},
\qquad
\Phi^{C_2}_{ij}:=\{\pm 2e_i,\pm 2e_j,\pm e_i\pm e_j\}\subset H_{ij}.
\]
特别地，对任意可见权
\[
\lambda\in \Lambda_B(X_6)\cup \Lambda_C(X_6)
\]
都满足
\[
\lambda_1\lambda_2\lambda_3=0.
\]
\end{theorem}
\begin{proof}
由定理 \ref{thm:window6-21-adjoint-weight-multiset}，\(B_3\) 规范下的 \(18\) 条非零权恰为
\[
\{\pm e_1,\pm e_2,\pm e_3\}
\sqcup
\{\pm e_i\pm e_j:\ 1\le i<j\le 3\},
\]
而 \(C_3\) 规范下的 \(18\) 条非零权恰为
\[
\{\pm 2e_1,\pm 2e_2,\pm 2e_3\}
\sqcup
\{\pm e_i\pm e_j:\ 1\le i<j\le 3\}.
\]
按零坐标的位置分组，在 \(B_3\) 规范下每个 \(H_{ij}\) 内恰出现
\[
\{\pm e_i,\pm e_j,\pm e_i\pm e_j\},
\]
即一个 \(B_2\) 根系；在 \(C_3\) 规范下每个 \(H_{ij}\) 内恰出现
\[
\{\pm 2e_i,\pm 2e_j,\pm e_i\pm e_j\},
\]
即一个 \(C_2\) 根系。
而 boundary 三词都被送到零向量，不影响非零支撑。
故所示并分解成立。

每个非零支撑点都落在某个坐标平面 \(H_{ij}\) 上，故至少有一个坐标为零，从而
\(
\lambda_1\lambda_2\lambda_3=0
\)
对所有可见权成立。证毕。
\end{proof}

\begin{proposition}[三轴选择律与指数特征函数的三阶湮灭]\label{prop:window6-b3c3-triaxial-selection-ideal-factorization}
定义可见矩泛函
\[
\mathcal M_B(P):=\sum_{w\in X_6}P(\Lambda_B(w)),
\qquad
\mathcal M_C(P):=\sum_{w\in X_6}P(\Lambda_C(w)),
\qquad
P\in \CC[x_1,x_2,x_3].
\]
则 \(\mathcal M_B\) 与 \(\mathcal M_C\) 都经由商代数
\[
\CC[x_1,x_2,x_3]/(x_1x_2x_3)
\]
因子化。
等价地，只要 \(a,b,c\ge 1\)，便有
\[
\sum_{w\in X_6}\Lambda_B(w)_1^a\Lambda_B(w)_2^b\Lambda_B(w)_3^c=0,
\qquad
\sum_{w\in X_6}\Lambda_C(w)_1^a\Lambda_C(w)_2^b\Lambda_C(w)_3^c=0.
\]
进一步，沿用定理 \ref{thm:window6-b3c3-adjoint-character-and-quartic-invariant} 的指数特征函数
\[
\Theta_B(t):=\sum_{w\in X_6}e^{\langle \Lambda_B(w),t\rangle},
\qquad
\Theta_C(t):=\sum_{w\in X_6}e^{\langle \Lambda_C(w),t\rangle},
\qquad
t=(t_1,t_2,t_3)\in \RR^3,
\]
则恒有
\[
\partial_{t_1}\partial_{t_2}\partial_{t_3}\Theta_B\equiv 0,
\qquad
\partial_{t_1}\partial_{t_2}\partial_{t_3}\Theta_C\equiv 0.
\]
\end{proposition}
\begin{proof}
由定理 \ref{thm:window6-b3c3-visible-support-three-levi-planes}，对每个可见权 \(\lambda\) 都有
\[
\lambda_1\lambda_2\lambda_3=0.
\]
因此任意落在理想 \((x_1x_2x_3)\) 中的多项式在全部支撑点上逐点为零，故 \(\mathcal M_B,\mathcal M_C\) 都通过
\(
\CC[x_1,x_2,x_3]/(x_1x_2x_3)
\)
因子化。
把该结论应用到单项式 \(x_1^a x_2^b x_3^c\) 即得幂矩消失式。

对指数特征函数，逐项求导可得
\[
\partial_{t_1}\partial_{t_2}\partial_{t_3}
e^{\langle \lambda,t\rangle}
=
\lambda_1\lambda_2\lambda_3\,e^{\langle \lambda,t\rangle}
=
0.
\]
对全部 \(\lambda=\Lambda_B(w)\) 或 \(\Lambda_C(w)\) 求和，即得所示恒等式。证毕。
\end{proof}

\begin{theorem}[quartic 缺陷的一维性]\label{thm:window6-b3c3-quartic-defect-onedim}
定义
\[
s_2(u):=u_1^2+u_2^2+u_3^2,
\qquad
s_4(u):=u_1^4+u_2^4+u_3^4,
\]
并记归一化八面体型调和四次式
\[
H_4^{\mathrm{oct}}(u):=s_4(u)-\frac35 s_2(u)^2.
\]
则对任意 \(u\in \RR^3\)，有精确恒等式
\[
\sum_{w\in X_6}\langle \Lambda_B(w),u\rangle^2=10\,s_2(u),
\qquad
\sum_{w\in X_6}\langle \Lambda_C(w),u\rangle^2=16\,s_2(u),
\]
以及
\[
\sum_{w\in X_6}\langle \Lambda_B(w),u\rangle^4
=
\frac{54}{5}s_2(u)^2-2\,H_4^{\mathrm{oct}}(u),
\]
\[
\sum_{w\in X_6}\langle \Lambda_C(w),u\rangle^4
=
\frac{144}{5}s_2(u)^2+28\,H_4^{\mathrm{oct}}(u).
\]
因此，window-\(6\) 可见部门在二阶上完全各向同性，而其四阶偏离只沿唯一的一维调和方向 \(\RR H_4^{\mathrm{oct}}\) 存活。
\end{theorem}
\begin{proof}
前两条二阶恒等式正是定理 \ref{thm:window6-b3c3-adjoint-second-moment-isotropy} 的内容。
又由定理 \ref{thm:window6-b3c3-adjoint-character-and-quartic-invariant}，
\[
\sum_{w\in X_6}\langle \Lambda_B(w),u\rangle^4
=
10s_4(u)+24\sum_{1\le i<j\le 3}u_i^2u_j^2,
\]
\[
\sum_{w\in X_6}\langle \Lambda_C(w),u\rangle^4
=
40s_4(u)+24\sum_{1\le i<j\le 3}u_i^2u_j^2.
\]
由
\[
s_2(u)^2=s_4(u)+2\sum_{1\le i<j\le 3}u_i^2u_j^2
\]
可化为
\[
\sum_{w\in X_6}\langle \Lambda_B(w),u\rangle^4
=
12s_2(u)^2-2s_4(u),
\]
\[
\sum_{w\in X_6}\langle \Lambda_C(w),u\rangle^4
=
12s_2(u)^2+28s_4(u).
\]
再用
\[
s_4(u)=H_4^{\mathrm{oct}}(u)+\frac35 s_2(u)^2
\]
代回，即得
\[
12s_2^2-2s_4
=
\frac{54}{5}s_2^2-2H_4^{\mathrm{oct}},
\]
\[
12s_2^2+28s_4
=
\frac{144}{5}s_2^2+28H_4^{\mathrm{oct}}.
\]
故四阶非径向偏离只沿 \(H_4^{\mathrm{oct}}\) 这一条调和方向出现。证毕。
\end{proof}

\begin{theorem}[quartic 极值方向的 axis--diagonal 反转]\label{thm:window6-b3c3-quartic-axis-diagonal-reversal}
定义
\[
M_{4,B}(u):=\sum_{w\in X_6}\langle \Lambda_B(w),u\rangle^4,
\qquad
M_{4,C}(u):=\sum_{w\in X_6}\langle \Lambda_C(w),u\rangle^4.
\]
若 \(|u|=1\)，则
\[
M_{4,B}(u)=12-2\sum_{i=1}^{3}u_i^4,
\qquad
M_{4,C}(u)=12+28\sum_{i=1}^{3}u_i^4.
\]
因此
\[
\min_{|u|=1}M_{4,B}(u)=10,
\qquad
\max_{|u|=1}M_{4,B}(u)=\frac{34}{3},
\]
\[
\min_{|u|=1}M_{4,C}(u)=\frac{64}{3},
\qquad
\max_{|u|=1}M_{4,C}(u)=40.
\]
并且极值方向精确为：
\begin{enumerate}
  \item \(B_3\) 规范下，最小值在坐标轴方向
  \[
  \{\pm e_1,\pm e_2,\pm e_3\}
  \]
  处取得，而最大值在体对角方向
  \[
  \frac{1}{\sqrt3}(\pm 1,\pm 1,\pm 1)
  \]
  处取得；
  \item \(C_3\) 规范下，上述极值方向完全反转。
\end{enumerate}
\end{theorem}
\begin{proof}
由定理 \ref{thm:window6-b3c3-quartic-defect-onedim}，在单位球面 \(|u|=1\) 上有
\[
H_4^{\mathrm{oct}}(u)=\sum_{i=1}^{3}u_i^4-\frac35.
\]
代回即可得到
\[
M_{4,B}(u)=12-2\sum_{i=1}^{3}u_i^4,
\qquad
M_{4,C}(u)=12+28\sum_{i=1}^{3}u_i^4.
\]

令
\[
\Sigma_4(u):=\sum_{i=1}^{3}u_i^4.
\]
在约束 \(|u|=1\) 下，由 Hölder 不等式或 Lagrange 乘子法，
\[
\frac13\le \Sigma_4(u)\le 1,
\]
其中 \(\Sigma_4(u)=1\) 当且仅当 \(u\) 为坐标轴方向，而 \(\Sigma_4(u)=1/3\) 当且仅当
\[
|u_1|=|u_2|=|u_3|=\frac{1}{\sqrt3}.
\]
把这两个极值代回上述线性公式，即得全部极值与极值方向。证毕。
\end{proof}

\begin{theorem}[B/C 规范差的全偶矩层级与 Laplace 支配]\label{thm:window6-b3c3-even-moment-gap-laplace-domination}
对任意整数 \(m\ge 1\) 与任意 \(u\in \RR^3\)，记
\[
M_{2m}^{B}(u):=\sum_{w\in X_6}\langle \Lambda_B(w),u\rangle^{2m},
\qquad
M_{2m}^{C}(u):=\sum_{w\in X_6}\langle \Lambda_C(w),u\rangle^{2m}.
\]
则有精确恒等式
\[
M_{2m}^{C}(u)-M_{2m}^{B}(u)
=
2\bigl(2^{2m}-1\bigr)\sum_{i=1}^{3}u_i^{2m}.
\]
等价地，指数特征函数满足
\[
\Theta_C(t)-\Theta_B(t)
=
2\sum_{i=1}^{3}\bigl(\cosh(2t_i)-\cosh(t_i)\bigr)\ge 0
\qquad
(t\in \RR^3),
\]
且当且仅当 \(t=0\) 时取等。
因此，对每个固定方向 \(u\)，标量随机变量 \(\langle \Lambda_C,u\rangle\) 在 Laplace 序上严格支配 \(\langle \Lambda_B,u\rangle\)。
\end{theorem}
\begin{proof}
由定理 \ref{thm:window6-21-adjoint-weight-multiset}，\(\Lambda_B\) 与 \(\Lambda_C\) 的差别只发生在六个轴向词上：\(B_3\) 规范给出
\[
\{\pm e_1,\pm e_2,\pm e_3\},
\]
而 \(C_3\) 规范把它们改写为
\[
\{\pm 2e_1,\pm 2e_2,\pm 2e_3\}.
\]
其余十二条双坐标根和三重零权完全相同。
因此
\[
M_{2m}^{C}(u)-M_{2m}^{B}(u)
=
\sum_{i=1}^{3}\Bigl((2u_i)^{2m}+(-2u_i)^{2m}-u_i^{2m}-(-u_i)^{2m}\Bigr),
\]
即得
\[
M_{2m}^{C}(u)-M_{2m}^{B}(u)
=
2\bigl(2^{2m}-1\bigr)\sum_{i=1}^{3}u_i^{2m}.
\]

对该恒等式按 \(m\) 求指数母函数，或直接利用定理 \ref{thm:window6-b3c3-adjoint-character-and-quartic-invariant} 中 \(\Theta_B,\Theta_C\) 的显式公式，便得到
\[
\Theta_C(t)-\Theta_B(t)
=
2\sum_{i=1}^{3}\bigl(\cosh(2t_i)-\cosh(t_i)\bigr).
\]
因为 \(\cosh(2x)>\cosh(x)\) 对每个 \(x\neq 0\) 都成立，故右端非负，且只有在 \(t_1=t_2=t_3=0\) 时为零。
于是 \(\Theta_C(t)\ge \Theta_B(t)\) 对所有 \(t\in \RR^3\) 成立，并在 \(t\neq 0\) 时严格大于，故 \(C_3\) 规范在每个方向上的 Laplace 变换都严格支配 \(B_3\) 规范。证毕。
\end{proof}

\begin{conjecture}[GFF 标度下可见部门的唯一 quartic 残差原理]\label{conj:window6-visible-gff-unique-quartic-residual}
设某个临界六顶点系综中，window-\(6\) 的可见线性观测在适当归一化后收敛到标量 Gaussian free field 的线性像。
则单点可见有效作用的局域展开应同时满足：
\begin{enumerate}
  \item 不存在真正的三轴 cubic 项；
  \item 二阶项必为标量刚度；
  \item 首个非高斯、非标量且可见的局域修正必与
  \[
  H_4^{\mathrm{oct}}(u)=u_1^4+u_2^4+u_3^4-\frac35|u|^4
  \]
  成比例；
  \item \(B_3\) 与 \(C_3\) 两种规范只改变这一 quartic 项的系数，不改变其形式。
\end{enumerate}
\end{conjecture}

\begin{remark}[唯一 quartic 残差原则的离散依据]\label{rem:window6-visible-gff-unique-quartic-residual-evidence}
命题 \ref{prop:window6-b3c3-triaxial-selection-ideal-factorization} 排除了全部真正的三轴单点耦合；定理 \ref{thm:window6-b3c3-quartic-defect-onedim} 表明二阶完全各向同性，而四阶缺陷空间仅为一维；定理 \ref{thm:window6-b3c3-even-moment-gap-laplace-domination} 又说明 \(B/C\) 规范差在所有偶矩层级上都只能通过
\[
\sum_{i=1}^{3}u_i^{2m}
\]
这一轴向纯幂族进入。
因此，若 window-\(6\) 的可见部门在连续极限中确有首个非高斯残差，则它既不应首先出现在 cubic 通道，也不应分裂成多个独立的张量方向；最先显现的可见局域异常只能是唯一的八面体型 quartic operator。
\end{remark}

\endinput


\begin{theorem}[阿贝尔化奇偶荷格的根--Cartan 三分裂与 boundary parity 的零权本质]\label{thm:window6-abelianized-parity-charge-root-cartan-splitting}
记
\[
\mathcal G_6^{\mathrm{bin}}:=\Aut(\Fold_6^{\mathrm{bin}}),
\]
并经由纤维符号映射 \(\mathrm{sgn}_6\) 的坐标识别写成
\[
\bigl(\mathcal G_6^{\mathrm{bin}}\bigr)^{\mathrm{ab}}\cong (\ZZ_2)^{X_6}.
\]
在 \(B_3\) 规范下令
\[
X_6^{\mathrm{short}}:=\{w\in X_6^{\mathrm{cyc}}:\ \mathrm{wt}(w)=1\},
\qquad
X_6^{\mathrm{long}}:=X_6^{\mathrm{cyc}}\setminus X_6^{\mathrm{short}}.
\]
并记相应的根集合为
\[
\Phi_{\mathrm{short}}:=\{\pm e_1,\pm e_2,\pm e_3\},
\qquad
\Phi_{\mathrm{long}}:=\{\pm e_i\pm e_j:\ 1\le i<j\le 3\}.
\]
则存在规范直和分解
\[
\bigl(\mathcal G_6^{\mathrm{bin}}\bigr)^{\mathrm{ab}}
\cong
(\ZZ_2)^{X_6^{\mathrm{short}}}
\oplus
(\ZZ_2)^{X_6^{\mathrm{long}}}
\oplus
(\ZZ_2)^{X_6^{\mathrm{bdry}}}
\cong
\FF_2[\Phi_{\mathrm{short}}]
\oplus
\FF_2[\Phi_{\mathrm{long}}]
\oplus
\mathfrak h_{\FF_2},
\]
其中
\[
\dim_{\FF_2}\FF_2[\Phi_{\mathrm{short}}]=6,\qquad
\dim_{\FF_2}\FF_2[\Phi_{\mathrm{long}}]=12,\qquad
\dim_{\FF_2}\mathfrak h_{\FF_2}=3.
\]
若改用 \(C_3\) 规范，则前两项仅交换长短根名称，而三项维数保持不变。

进一步，记
\[
Z_{\partial}:=\langle \tau_w:\ w\in X_6^{\mathrm{bdry}}\rangle\cong (\ZZ_2)^3.
\]
则 Abel 化映射把 \(Z_{\partial}\) 同构到第三项 \(\mathfrak h_{\FF_2}\) 上。
因此 boundary sheet parity 在该伴随权模型中对应的正是零权的 Cartan 部分，而不是根方向电荷。
\end{theorem}
\begin{proof}
由定理 \ref{thm:window6-cyclic-weight-threshold-root-length} 与定理 \ref{thm:window6-21-adjoint-weight-multiset}，集合 \(X_6\) 已被规范分解为
\[
X_6=X_6^{\mathrm{short}}\sqcup X_6^{\mathrm{long}}\sqcup X_6^{\mathrm{bdry}},
\qquad
|X_6^{\mathrm{short}}|=6,\quad |X_6^{\mathrm{long}}|=12,\quad |X_6^{\mathrm{bdry}}|=3.
\]
另一方面，定理 \ref{thm:window6-foldbin-gauge-abelianization-even-parity} 给出
\[
\bigl(\mathcal G_6^{\mathrm{bin}}\bigr)^{\mathrm{ab}}\cong (\ZZ_2)^{21},
\]
并且该 Abel 化通过 \(\mathrm{sgn}_6\) 被识别为以 \(X_6\) 为坐标集的奇偶荷格。
因此对上述三块分割逐块取坐标子格，便得到所示的规范直和分解。

再由推论 \ref{cor:fold-bin-boundary-center-z2}，boundary 三条二点纤维的纤维内换位 involution 生成一个规范中心子群
\(
Z_{\partial}\cong(\ZZ_2)^3
\)，其在 Abel 化中的像恰对应 boundary 三个坐标方向。
这正是第三项 \((\ZZ_2)^{X_6^{\mathrm{bdry}}}=\mathfrak h_{\FF_2}\)。
故 boundary parity 的本质是零权的 Cartan 部分，而非根空间部分。证毕。
\end{proof}

\begin{theorem}[window-\(6\) 的 \texorpdfstring{$C_6$}{C6} 轨道分解与 \texorpdfstring{$B_3$}{B3} 长短根分区的交会刚性]\label{thm:window6-c6-orbit-b3-root-intersection-rigidity}
沿用 \(B_3\) 规范下的分块
\[
X_6^{\mathrm{cyc}}=X_6^{\mathrm{short}}\sqcup X_6^{\mathrm{long}}.
\]
则表 \ref{tab:window6_c6_orbit_decomposition} 给出的五条 \(C_6\) 轨道中，恰有且仅有一条长度为 \(6\) 的轨道满足
\[
X_6^{\mathrm{short}}
=
\{\texttt{000001},\texttt{000010},\texttt{000100},\texttt{001000},\texttt{010000},\texttt{100000}\},
\]
其余四条轨道全部包含于 \(X_6^{\mathrm{long}}\)。
特别地，另一条长度为 \(6\) 的轨道
\[
\{\texttt{000101},\texttt{001010},\texttt{010001},\texttt{010100},\texttt{100010},\texttt{101000}\}
\]
整体落在长根扇区之内。

从而，同时满足“在每条 \(C_6\) 轨道上常值”与“Weyl-不变”的函数空间恰为
\[
\left\{
f:X_6^{\mathrm{cyc}}\to\CC:\ 
f\text{ 在每条 }C_6\text{ 轨道上常值且为 Weyl-不变}
\right\}
=
\operatorname{span}_{\CC}\{\mathbf 1_{X_6^{\mathrm{short}}},\mathbf 1_{X_6^{\mathrm{long}}}\},
\]
因而其维数恰为 \(2\)。
\end{theorem}

\begin{proof}
由表 \ref{tab:window6_c6_orbit_decomposition}，\(X_6^{\mathrm{cyc}}\) 的两条长度为 \(6\) 的轨道分别为
\[
\{\texttt{000001},\texttt{000010},\texttt{000100},\texttt{001000},\texttt{010000},\texttt{100000}\}
\]
与
\[
\{\texttt{000101},\texttt{001010},\texttt{010001},\texttt{010100},\texttt{100010},\texttt{101000}\}.
\]
其中第一条轨道恰由全部 \(\mathrm{wt}=1\) 的六个词组成，故由定理
\ref{thm:window6-cyclic-weight-threshold-root-length} 正好等于 \(X_6^{\mathrm{short}}\)。
第二条长度为 \(6\) 的轨道在表 \ref{tab:fold6_b3c3_root_dictionary_b3} 中逐项对应长根，因而整体落在 \(X_6^{\mathrm{long}}\)。

其余三条轨道分别由 \(\texttt{000000}\)、\(\{\texttt{010101},\texttt{101010}\}\) 与
\(\{\texttt{001001},\texttt{010010},\texttt{100100}\}\) 生成；
它们的 Hamming 权分别为 \(0\)、\(3\) 与 \(2\)，故仍由定理
\ref{thm:window6-cyclic-weight-threshold-root-length} 全部落在 \(X_6^{\mathrm{long}}\)。
因此五条 \(C_6\) 轨道与长短根二分块的交会类型被唯一确定为：
唯一一条长度为 \(6\) 的轨道等于短根扇区，其余四条轨道全为长根扇区的子块。

现在设 \(f:X_6^{\mathrm{cyc}}\to\CC\) 在每条 \(C_6\) 轨道上常值且 Weyl-不变。
由于 \(X_6^{\mathrm{short}}\) 本身就是一条单独的 \(C_6\) 轨道，\(f\) 在其上只能取一个常数。
另一方面，\(X_6^{\mathrm{long}}\) 虽由四条 \(C_6\) 轨道组成，但 Weyl-不变性强迫 \(f\) 在整块 \(X_6^{\mathrm{long}}\) 上取同一个常数。
故该函数空间恰由 \(\mathbf 1_{X_6^{\mathrm{short}}}\) 与 \(\mathbf 1_{X_6^{\mathrm{long}}}\) 张成，维数为 \(2\)。证毕。
\end{proof}

\paragraph{visible Cartan quotient、syzygy lattice 与二元码闭链}
在 \(B_3\) 可见根字典、boundary 三基、四次调和检测器与 visible Laplace 支配都已固定之后，window-\(6\) 的 \(21\) 个可见模态还可继续压缩为一条更刚性的整数线性链：
boundary 三词先给出一个规范的 rank-\(3\) Cartan quotient，其核形成一个秩 \(18\) 的 syzygy lattice；该 lattice 的 Gram 判别式与 visible matroid 的全部三阶体积预算同时锁定到 \(11^3\)；四阶层面恰有唯一的 boundary counterterm 能把 cyclic cloud 各向同性化，而六阶则出现不可消去的各向异性障碍；最后，这个模 \(2\) Cartan quotient 自然生成一对互补的 \([21,3,9]\) 与 \([21,18,2]\) 码。

\leanverified{paper\_window6\_visible\_cartan\_quotient\_syzygy\_splitting}
\begin{theorem}[visible Cartan quotient 的规范分裂与 syzygy 基]\label{thm:window6-visible-cartan-quotient-syzygy-splitting}
固定 boundary 排序
\[
b_1:=\texttt{100001},
\qquad
b_2:=\texttt{100101},
\qquad
b_3:=\texttt{101001},
\]
并记表 \ref{tab:fold6_b3c3_root_dictionary_b3} 的 \(B_3\) 根字典为
\[
\alpha:X_6^{\mathrm{cyc}}\longrightarrow \Phi(B_3)\subset \ZZ^3.
\]
定义可见向量
\[
v_w:=
\begin{cases}
\alpha(w),& w\in X_6^{\mathrm{cyc}},\\
e_i,& w=b_i\in X_6^{\mathrm{bdry}},
\end{cases}
\]
以及整数线性映射
\[
\Pi_6:\ZZ^{X_6}\longrightarrow \ZZ^3,
\qquad
\Pi_6(\delta_w):=v_w,
\]
其中 \(\{\delta_w\}_{w\in X_6}\) 为标准坐标基。
则存在规范分裂正合列
\[
0\longrightarrow \Lambda_6^{\mathrm{syz}}
\longrightarrow \ZZ^{X_6}
\xrightarrow{\ \Pi_6\ }
\ZZ^3
\longrightarrow 0,
\]
其分裂由
\[
s:\ZZ^3\to \ZZ^{X_6},
\qquad
s(e_i)=\delta_{b_i}
\]
给出。更精确地，
\[
\Lambda_6^{\mathrm{syz}}=\ker \Pi_6
\]
是秩 \(18\) 的自由格，并有显式基
\[
r_w:=\delta_w-\sum_{i=1}^{3}\alpha_i(w)\,\delta_{b_i},
\qquad
w\in X_6^{\mathrm{cyc}}.
\]
因此
\[
\ZZ^{X_6}\cong \ZZ^3\oplus \ZZ^{18},
\]
其中三维直和因子是 visible Cartan 层，十八维直和因子是 cyclic syzygy 层。
\end{theorem}
\begin{proof}
由于 \(v_{b_i}=e_i\)，映射 \(\Pi_6\) 显然满射，且
\[
\Pi_6\circ s=\mathrm{id}_{\ZZ^3},
\]
故该正合列分裂。

对每个 \(w\in X_6^{\mathrm{cyc}}\)，有
\[
\Pi_6(r_w)
=
\alpha(w)-\sum_{i=1}^{3}\alpha_i(w)e_i
=
0,
\]
故 \(r_w\in \ker \Pi_6\)。
这些向量线性无关，因为每个 \(r_w\) 在其专属坐标 \(\delta_w\) 上系数为 \(1\)，而其余 \(r_{w'}\) 在该坐标上系数为 \(0\)。

另一方面，
\[
\operatorname{rank}\ker \Pi_6
=
|X_6|-3
=
21-3
=
18
=
|X_6^{\mathrm{cyc}}|,
\]
故 \(\{r_w\}_{w\in X_6^{\mathrm{cyc}}}\) 即为 \(\ker \Pi_6\) 的一组基。

最后，把该构造模 \(2\) 约化，并结合定理 \ref{thm:window6-foldbin-gauge-abelianization-even-parity} 中
\[
\bigl(\mathcal G_6^{\mathrm{bin}}\bigr)^{\mathrm{ab}}\cong (\ZZ_2)^{21}
\]
及推论 \ref{cor:fold-bin-boundary-center-z2} 给出的 boundary 三维坐标子格，可得 anomaly lattice 上的同样分裂。证毕。
\end{proof}

\begin{theorem}[syzygy lattice 的 Gram 谱与判别式]\label{thm:window6-syzygy-gram-spectrum-discriminant}
\leanverified{paper\_window6\_syzygy\_gram\_spectrum\_discriminant}
令
\[
R_6:=\bigl[r_w\bigr]_{w\in X_6^{\mathrm{cyc}}}\in M_{21\times 18}(\ZZ)
\]
为定理 \ref{thm:window6-visible-cartan-quotient-syzygy-splitting} 的 syzygy 基矩阵，并记
\[
K_6:=R_6^{\mathsf T}R_6.
\]
则
\[
\operatorname{Spec}(K_6)=\{11^{(3)},1^{(15)}\},
\]
从而
\[
\det K_6=11^3,
\qquad
\operatorname{tr}K_6=48.
\]
特别地，
\[
\operatorname{disc}(\Lambda_6^{\mathrm{syz}})=11^3.
\]
因此 \(\Lambda_6^{\mathrm{syz}}\) 并非均匀刚性，而是精确裂解为 \(15\) 个单位刚度方向与 \(3\) 个刚度 \(11\) 的特异方向；后者的个数恰与 Cartan 秩相同。
\end{theorem}
\begin{proof}
记
\[
A_6^{\mathrm{cyc}}
:=
\bigl[\alpha(w)\bigr]_{w\in X_6^{\mathrm{cyc}}}
\in M_{3\times 18}(\ZZ).
\]
由
\[
r_w=\delta_w-\sum_{i=1}^{3}\alpha_i(w)\delta_{b_i}
\]
的定义直接可见
\[
K_6=I_{18}+\bigl(A_6^{\mathrm{cyc}}\bigr)^{\mathsf T}A_6^{\mathrm{cyc}}.
\]
另一方面，定理 \ref{thm:window6-b3c3-adjoint-second-moment-isotropy} 给出 cyclic 根云的二阶矩恒等式
\[
A_6^{\mathrm{cyc}}\bigl(A_6^{\mathrm{cyc}}\bigr)^{\mathsf T}
=
\sum_{w\in X_6^{\mathrm{cyc}}}\alpha(w)\alpha(w)^{\mathsf T}
=
10I_3.
\]
故 \(\bigl(A_6^{\mathrm{cyc}}\bigr)^{\mathsf T}A_6^{\mathrm{cyc}}\) 的非零特征值恰为 \(10\)（重数 \(3\)），其余 \(15\) 个特征值为 \(0\)。
对其加上 \(I_{18}\) 后便得到
\[
\operatorname{Spec}(K_6)=\{11^{(3)},1^{(15)}\}.
\]
于是
\[
\det K_6=11^3,
\qquad
\operatorname{tr}K_6=3\cdot 11+15\cdot 1=48.
\]
由于 \(\{r_w\}\) 是 \(\Lambda_6^{\mathrm{syz}}\) 的整基，\(\det K_6\) 正是该格的判别式。证毕。
\end{proof}

\begin{theorem}[quartic isotropization 的唯一 boundary counterterm]\label{thm:window6-visible-quartic-shell-isotropization}
\leanverified{paper\_window6\_visible\_quartic\_shell\_isotropization}
对 \(d\in\{2,3,4\}\)，定义 cyclic 壳层
\[
\Phi_{6,d}^{\mathrm{cyc}}
:=
\{\alpha(w):\ w\in X_6^{\mathrm{cyc}},\ d_6^{\mathrm{bin}}(w)=d\}.
\]
对参数 \((a_2,a_3,a_4,b)\in \CC^4\)，定义四次型
\[
\mathcal Q_{a_2,a_3,a_4,b}(x):=
a_2\sum_{\beta\in \Phi_{6,2}^{\mathrm{cyc}}}\langle x,\beta\rangle^4
+
a_3\sum_{\beta\in \Phi_{6,3}^{\mathrm{cyc}}}\langle x,\beta\rangle^4
+
a_4\sum_{\beta\in \Phi_{6,4}^{\mathrm{cyc}}}\langle x,\beta\rangle^4
+
b\sum_{i=1}^{3}x_i^4.
\]
则存在常数 \(c\) 使
\[
\mathcal Q_{a_2,a_3,a_4,b}(x)=c|x|^4
\]
当且仅当
\[
a_2=a_3=a_4=:a,
\qquad
b=2a,
\qquad
c=12a.
\]
从而得到唯一的 shell-constant 正权公式
\[
\sum_{w\in X_6^{\mathrm{cyc}}}\langle x,\alpha(w)\rangle^4
+
2\sum_{i=1}^{3}x_i^4
=
12|x|^4.
\]
因此 boundary 三重态在四阶通道中不是附加饰项，而是把 cyclic root cloud 恢复为精确 \(O(3)\)-isotropy 的唯一 counterterm。
\end{theorem}
\begin{proof}
由 window-\(6\) 的刚性直方图 \(\#\{d_6=2\}=8\)、\(\#\{d_6=3\}=4\)、\(\#\{d_6=4\}=9\) 可知，cyclic 壳层的基数分别为 \(5,4,9\)，其中 \(d_6=2\) 的另外 \(3\) 个状态正好是 boundary 三词。
结合表 \ref{tab:fold6_b3c3_root_dictionary_b3} 的根字典，直接展开可得
\[
\sum_{\beta\in\Phi_{6,2}^{\mathrm{cyc}}}\langle x,\beta\rangle^4
=
x_1^4+4x_2^4+4x_3^4+24x_2^2x_3^2,
\]
\[
\sum_{\beta\in\Phi_{6,3}^{\mathrm{cyc}}}\langle x,\beta\rangle^4
=
4x_1^4+4x_3^4+24x_1^2x_3^2,
\]
\[
\sum_{\beta\in\Phi_{6,4}^{\mathrm{cyc}}}\langle x,\beta\rangle^4
=
5x_1^4+6x_2^4+2x_3^4+24x_1^2x_2^2.
\]
若
\[
\mathcal Q_{a_2,a_3,a_4,b}(x)=c(x_1^2+x_2^2+x_3^2)^2,
\]
则比较交叉项
\[
x_1^2x_2^2,\qquad x_1^2x_3^2,\qquad x_2^2x_3^2
\]
的系数，得到
\[
24a_4=2c,\qquad 24a_3=2c,\qquad 24a_2=2c,
\]
故
\[
a_2=a_3=a_4=\frac{c}{12}.
\]
再比较 \(x_1^4\) 系数，得到
\[
\frac{c}{12}+4\cdot\frac{c}{12}+5\cdot\frac{c}{12}+b=c,
\]
即
\[
b=\frac{c}{6}=2a.
\]
反之，把 \(a_2=a_3=a_4=a\) 与 \(b=2a\) 代回，直接化简便得
\[
\mathcal Q_{a,a,a,2a}(x)=12a\,|x|^4.
\]
证毕。
\end{proof}

\begin{proposition}[六阶旋转不变的不可实现性]\label{prop:window6-visible-degree6-isotropy-obstruction}
设 \(\{c_w\}_{w\in X_6}\subset \CC\) 任意，并定义六次齐次多项式
\[
P_c(x):=\sum_{w\in X_6}c_w\langle x,v_w\rangle^6.
\]
则 \(P_c\) 中单项式 \(x_1^2x_2^2x_3^2\) 的系数恒等于零。因而不存在非零常数 \(\lambda\) 使
\[
P_c(x)=\lambda |x|^6.
\]
特别地，window-\(6\) 的 visible support 上精确旋转不变的偶矩只能维持到四阶；六阶处即出现不可消去的各向异性障碍。
\end{proposition}
\leanverified{paper\_window6\_visible\_degree6\_isotropy\_obstruction}
\begin{proof}
由 \(v_w\) 的定义，任意可见向量都至少有一个零坐标：cyclic 根向量只支撑在至多两条坐标轴上，而 boundary 三向量正是 \(e_1,e_2,e_3\)。
写
\[
v_w=(a_1,a_2,a_3),
\]
则
\[
\langle x,v_w\rangle^6=(a_1x_1+a_2x_2+a_3x_3)^6
\]
中，单项式 \(x_1^2x_2^2x_3^2\) 的系数为
\[
90\,a_1^2a_2^2a_3^2,
\]
而这在某个 \(a_i=0\) 时恒为零。
故对所有 \(w\) 的线性组合 \(P_c\)，该单项式系数始终为零。

另一方面，
\[
|x|^6=(x_1^2+x_2^2+x_3^2)^3
\]
中 \(x_1^2x_2^2x_3^2\) 的系数等于 \(6\)。
因此若 \(P_c(x)=\lambda |x|^6\)，必有 \(\lambda=0\)。证毕。
\end{proof}

\begin{theorem}[visible matroid 的 minor census 与 cubic volume budget]\label{thm:window6-visible-matroid-minor-census}
令
\[
A_6:=\bigl[v_w\bigr]_{w\in X_6}\in M_{3\times 21}(\ZZ).
\]
对每个三元子集 \(I\subset X_6\)（\(|I|=3\)），记 \(A_6[I]\) 为对应的 \(3\times 3\) 列子矩阵。则
\[
\det A_6[I]\in\{0,\pm 1,\pm 2\}
\qquad
(\forall I\in \tbinom{X_6}{3}).
\]
若定义
\[
N_r:=
\#\Bigl\{
I\in\tbinom{X_6}{3}:\ |\det A_6[I]|=r
\Bigr\},
\qquad
r\in\{0,1,2\},
\]
则有精确计数
\[
N_1=675,
\qquad
N_2=164,
\qquad
N_0=491.
\]
特别地，
\[
\sum_{I\in\binom{X_6}{3}}\det\!\bigl(A_6[I]\bigr)^2
=
675+4\cdot 164
=
1331
=
11^3.
\]
因此，window-\(6\) 的 visible matroid 并非一般的 rank-\(3\) 整矩阵结构，而是被一个 cubic budget \(11^3\) 精确控制。
\end{theorem}
\begin{proof}
每一列 \(v_w\) 都属于 \(\ZZ^3\)，且列向量的平方范数至多为 \(2\)。
由 Hadamard 不等式，
\[
|\det A_6[I]|^2\le 2\cdot 2\cdot 2=8,
\]
故 \(|\det A_6[I]|\le 2\)。又行列式是整数，于是只能取 \(0,\pm 1,\pm 2\)。

下面计算 \(N_1\)。把 \(A_6\) 模 \(2\) 约化为
\[
\bar A_6\in M_{3\times 21}(\FF_2).
\]
由表 \ref{tab:fold6_b3c3_root_dictionary_b3} 与 boundary 三基可知，\(\bar A_6\) 的列只落在六种非零类型上：
\[
100,\ 010,\ 001
\quad\text{各重数 }3;
\qquad
110,\ 101,\ 011
\quad\text{各重数 }4.
\]
三列子矩阵有奇行列式，当且仅当这三列在 \(\FF_2^3\) 中线性无关。按类型分三类计数：

第一类，三个轴型 \(100,010,001\) 构成唯一一组基，贡献
\[
3\cdot 3\cdot 3=27.
\]
第二类，两个不同轴型加一个双坐标型。先选两条轴，有 \(\binom 32=3\) 种；与其和相等的那一条双坐标型会造成线性相关，故可选双坐标型恰有 \(2\) 个，因此共有 \(6\) 组类型基，每组贡献
\[
3\cdot 3\cdot 4=36,
\]
总计 \(216\)。
第三类，一条轴型加两条不同双坐标型。任取轴型有 \(3\) 种，再从三条双坐标型中任选两条有 \(\binom 32=3\) 种；这 \(9\) 组类型全都线性无关，每组贡献
\[
3\cdot 4\cdot 4=48,
\]
总计 \(432\)。
故
\[
N_1=27+216+432=675.
\]

另一方面，由定理 \ref{thm:window6-b3c3-adjoint-second-moment-isotropy}，cyclic 部分给出 \(10I_3\)；再加上 boundary 三基的贡献
\[
e_1e_1^{\mathsf T}+e_2e_2^{\mathsf T}+e_3e_3^{\mathsf T}=I_3,
\]
得到
\[
A_6A_6^{\mathsf T}=11I_3.
\]
故由 Cauchy--Binet，
\[
\sum_{I\in\binom{X_6}{3}}\det\!\bigl(A_6[I]\bigr)^2
=
\det(A_6A_6^{\mathsf T})
=
\det(11I_3)
=
11^3
=
1331.
\]
因此
\[
N_1+4N_2=1331,
\]
代入 \(N_1=675\) 可得
\[
N_2=\frac{1331-675}{4}=164.
\]
最后
\[
N_0=\binom{21}{3}-N_1-N_2
=
1330-675-164
=
491.
\]
证毕。
\end{proof}

\leanverified{paper\_window6\_visible\_matroid\_minor\_census}
\leanverified{paper\_window6\_cartan\_kernel\_code\_parameters}
\begin{theorem}[Cartan kernel code 与其对偶码的精确参数]\label{thm:window6-cartan-kernel-code-parameters}
令 \(\bar A_6\in M_{3\times 21}(\FF_2)\) 为 \(A_6\) 的模 \(2\) 约化，并定义
\[
C_6^\perp:=\operatorname{rowspan}_{\FF_2}(\bar A_6)\subset \FF_2^{21},
\qquad
C_6:=\ker \bar A_6.
\]
则：
\begin{enumerate}
  \item \(C_6^\perp\) 是一个 \([21,3,9]\) 码，其 Hamming 重量枚举多项式恰为
  \[
  W_{C_6^\perp}(z)=1+z^9+3z^{11}+3z^{14}.
  \]
  \item \(C_6\) 是一个 \([21,18,2]\) 码，且其 Hamming 重量枚举由 MacWilliams 变换写成
  \[
  W_{C_6}(x,y)
  =
  \frac18\Bigl(
  (x+y)^{21}
  +(x+y)^{12}(x-y)^9
  +3(x+y)^{10}(x-y)^{11}
  +3(x+y)^7(x-y)^{14}
  \Bigr).
  \]
\end{enumerate}
因此，window-\(6\) 的 \(21\) 比特 anomaly lattice 在 mod-\(2\) Cartan 投影之下自然生成一对互为正交补的二元线性码：三维对偶码记录 visible Cartan 信息，而十八维 kernel code 记录全部 gauge-neutral 组合。
\end{theorem}
\begin{proof}
由 \(\bar A_6\) 中三条 boundary 列正是标准基 \(e_1,e_2,e_3\)，可知其行秩为 \(3\)，故
\[
\dim C_6^\perp=3,
\qquad
\dim C_6=21-3=18.
\]

再由定理 \ref{thm:window6-visible-matroid-minor-census} 证明中给出的六类二进制列型，可直接计算这 \(8\) 个行空间码字的重量：
零码字重量为 \(0\)；三条单独的行各有重量
\[
3+4+4=11;
\]
任意两行之和的重量为 \(14\)；三行总和的重量为 \(9\)。
于是
\[
W_{C_6^\perp}(z)=1+z^9+3z^{11}+3z^{14},
\]
其最小非零重量为 \(9\)，故 \(C_6^\perp\) 的参数为 \([21,3,9]\)。

对偶码 \(C_6\) 的重量枚举由 MacWilliams 恒等式给出：
\[
W_{C_6}(x,y)
=
\frac1{|C_6^\perp|}
W_{C_6^\perp}(x+y,x-y),
\]
而 \(|C_6^\perp|=2^3=8\)，代入上式即得所示闭式。

最后，为证明 \(d(C_6)=2\)，只需注意
\[
\bar A_6(\delta_{\texttt{100000}}+\delta_{\texttt{100001}})=0,
\]
因为 \(\texttt{100000}\) 与 \(b_1=\texttt{100001}\) 在 \(\FF_2^3\) 中都映为 \(100\)。
故 \(C_6\) 含有重量 \(2\) 的非零码字，从而其最小距离恰为 \(2\)。证毕。
\end{proof}

\noindent
由此，window-\(6\) 的可见 \(21\) 点模型不再只是一个根字典或纤维直方图，而被压缩为同一条整数线性闭链：
boundary 三基给出 visible Cartan quotient，
其核形成判别式 \(11^3\) 的秩 \(18\) syzygy lattice；
四阶处存在且仅存在一个 boundary counterterm 能把 cyclic cloud 恢复为精确 \(O(3)\)-isotropy，
而六阶处又出现不可消去的各向异性；
最后，这个 quotient 的模 \(2\) 影子恰生成一对 \([21,3,9]\) 与 \([21,18,2]\) 的正交补码。
于是，Cartan quotient、syzygy lattice、quartic counterterm 与 binary code 四层结构便在同一对象上闭合。



